\chapter{Ricci Curvature Comparison}
\label{chap:pet-ricci-curvature-comparison}

\newcommand{\Ric}{\operatorname{Ric}}
\newcommand{\Hess}{\operatorname{Hess}}
\newcommand{\vol}{\operatorname{vol}}
\newcommand{\diam}{\operatorname{diam}}
\newcommand{\Iso}{\mathrm{Iso}}
\newcommand{\snk}{\mathrm{sn}}

Faithful transcription of Petersen, \emph{Riemannian Geometry} (GTM 171, 3rd ed.),
Chapter 7. The chapter develops the analytic consequences of a lower Ricci
curvature bound: the trace identities for the Laplacian and volume form of a
distance function and the resulting Bishop--Gromov absolute and relative
volume comparison theorems (§7.1.1--7.1.2); the maximum principle for
functions with one-sided Hessian or Laplacian bounds in the barrier
(support-function) sense, the Laplacian comparison theorem of Calabi, the
segment inequality of Cheeger and Colding and the weak Poincaré and
Poincaré--Sobolev inequalities it implies, and the Rellich compactness theorem
(§7.1.3--7.1.5); applications to finiteness of fundamental groups and to the
rigidity of manifolds attaining the maximal diameter allowed by Myers'
theorem (§7.2); the Cheeger--Gromoll splitting theorem and the resulting
structure theory for compact manifolds of nonnegative Ricci curvature (§7.3);
a guide to further reading (§7.4); and the chapter's exercises (§7.5).

\section{Volume Comparison}

\subsection{The Fundamental Equations}

Throughout §7.1, $(M,g)$ is a complete Riemannian $n$-manifold and
$r(x)=|xp|$ a distance function, smooth on an open set $U\subset M$ (from
§7.1.3 on, $U$ is the image of the interior of the segment domain at $p$, see
§5.7.3). Write $\partial_r=\nabla r$. Recall the fundamental equations for the
metric along $r$ (Proposition 3.2.11):
\[
  L_{\partial_r}g=2\Hess r,
  \qquad
  (\nabla_{\partial_r}\Hess r)(X,Y)+\Hess^2r(X,Y)=-R(X,\partial_r,\partial_r,Y),
\]
where $\Hess r(X,Y)=g(\nabla_X\partial_r,Y)$ and
$\Hess^2r(X,Y)=g(\nabla_X\partial_r,\nabla_Y\partial_r)$ (writing $\Hess r$
also for the associated self-adjoint operator). There is an analogous pair of
equations for the volume form.

\begin{proposition}[Proposition 7.1.1]
  \label{prop:pet-ch7-volume-form-laplacian-trace}
  \lean{PetersenLib.volumeFormLaplacianTraceIdentities}
  The volume form $\vol$ and Laplacian $\Delta r$ of a smooth distance
  function $r$ satisfy
  \[
    \text{(tr1)}\quad L_{\partial_r}\vol=\Delta r\,\vol,
    \qquad
    \text{(tr2)}\quad \partial_r\Delta r+\frac{(\Delta r)^2}{n-1}
    \le\partial_r\Delta r+|\Hess r|^2=-\Ric(\partial_r,\partial_r).
  \]
  \source{petersen:page-0291,petersen:page-0292}
\end{proposition}
\begin{proof}
  \uses{prop:pet-ch7-volume-form-laplacian-trace}
  (tr1) is one of the equivalent characterizations of $\Delta r$ established
  in §2.1.3, recalled here for use in this chapter.

  For (tr2), fix $p\in U$ and choose an orthonormal frame $E_1,\dots,E_n$ near
  $p$ that is parallel along the integral curve of $\partial_r$ through $p$
  (i.e.\ $\nabla_{\partial_r}E_i=0$); set $X=Y=E_i$ in the second fundamental
  equation and sum over $i$. On the right,
  \[
    \sum_{i=1}^nR(E_i,\partial_r,\partial_r,E_i)=\Ric(\partial_r,\partial_r).
  \]
  On the left, parallelism of the frame lets the derivative pass through the
  trace,
  \[
    \sum_{i=1}^n(\nabla_{\partial_r}\Hess r)(E_i,E_i)
    =\partial_r\sum_{i=1}^n\Hess r(E_i,E_i)=\partial_r\Delta r,
  \]
  while, using $\Hess r(X,Y)=g(\nabla_X\partial_r,Y)$ and expanding
  $\nabla_{E_i}\partial_r$ in the frame,
  \[
    \sum_{i=1}^n\Hess^2r(E_i,E_i)
    =\sum_{i=1}^ng(\nabla_{E_i}\partial_r,\nabla_{E_i}\partial_r)
    =\sum_{i,j=1}^ng(\nabla_{E_i}\partial_r,E_j)^2
    =|\Hess r|^2.
  \]
  Together these give the equality
  $\partial_r\Delta r+|\Hess r|^2=-\Ric(\partial_r,\partial_r)$.

  It remains to show $(\Delta r)^2/(n-1)\le|\Hess r|^2$. Take
  $E_1=\partial_r$; since $|\nabla r|\equiv1$ forces $\Hess
  r(\partial_r,\cdot)=0$ (differentiate $g(\partial_r,\partial_r)=1$ along any
  direction), the sum defining $|\Hess r|^2$ has no $i=1$ or $j=1$ terms:
  \[
    |\Hess r|^2=\sum_{i,j=1}^ng(\nabla_{E_i}\partial_r,E_j)^2
    =\sum_{i,j=2}^ng(\nabla_{E_i}\partial_r,E_j)^2.
  \]
  Writing $A=[g(\nabla_{E_i}\partial_r,E_j)]_{2\le i,j\le n}$, an
  $(n-1)\times(n-1)$ matrix, the Cauchy--Schwarz inequality applied to the
  Frobenius inner product of $A$ against the identity matrix $I_{n-1}$ gives
  \[
    (\operatorname{tr}A)^2=\langle A,I_{n-1}\rangle^2
    \le|A|^2|I_{n-1}|^2=(n-1)|A|^2,
  \]
  i.e.\ $|A|^2\ge(\operatorname{tr}A)^2/(n-1)$. Since
  $\operatorname{tr}A=\sum_{i=2}^ng(\nabla_{E_i}\partial_r,E_i)=\Delta r$ (the
  $i=1$ term of the full trace vanishes, as just noted) and $|A|^2=|\Hess
  r|^2$, this is exactly $|\Hess r|^2\ge(\Delta r)^2/(n-1)$, completing the
  proof of (tr2).
\end{proof}

Work in exponential polar coordinates at $p$: write $g=dr^2+g_r$ and let
$\vol_{n-1}$ denote the canonical volume form on $S^{n-1}(1)$. The volume form
of $g$ then factors as $\vol=\lambda(r,\theta)\,dr\wedge\vol_{n-1}$,
$\theta\in S^{n-1}$, for a positive function $\lambda$, the \emph{polar
density}. Since $L_{\partial_r}dr=0$ and $L_{\partial_r}\vol_{n-1}=0$ (both
are pulled back from the unit-sphere factor, independent of $r$), applying
(tr1) to this expression gives $L_{\partial_r}\vol=\partial_r(\lambda)\,dr
\wedge\vol_{n-1}$, so (tr1) reduces to the scalar ODE
$\partial_r\lambda=\lambda\,\Delta r$.

\begin{definition}[polar density and its constant-curvature model]
  \label{def:pet-ch7-polar-density}
  \lean{PetersenLib.polarDensity,PetersenLib.polarDensityConstantCurvature}
  \uses{prop:pet-ch7-volume-form-laplacian-trace}
  In exponential polar coordinates $g=dr^2+g_r$ at $p$, the \emph{polar
  density} $\lambda(r,\theta)>0$ is defined by
  $\vol=\lambda\,dr\wedge\vol_{n-1}$; it satisfies
  $\partial_r\lambda=\lambda\,\Delta r$. In the constant curvature $k$ space
  form, $g_k=dr^2+\snk_k^2(r)\,ds_{n-1}^2$, and the corresponding density
  $\lambda_k(r)=\snk_k^{n-1}(r)$ satisfies the same relation with
  $\Delta r=(n-1)\snk_k'(r)/\snk_k(r)$:
  \[
    \partial_r\big(\snk_k^{n-1}(r)\big)
    =(n-1)\frac{\snk_k'(r)}{\snk_k(r)}\,\snk_k^{n-1}(r).
  \]
  \source{petersen:page-0292,petersen:page-0293}
\end{definition}

\subsection{Volume Estimation}

With \cref{prop:pet-ch7-volume-form-laplacian-trace} and the polar density in
hand we obtain volume estimates analogous to the metric and Hessian
comparison estimates of §6.4, now under a lower Ricci (rather than sectional)
curvature bound.

\begin{lemma}[Lemma 7.1.2 — Ricci Comparison]
  \label{lem:pet-ch7-ricci-comparison-mean-curvature}
  \lean{PetersenLib.ricciMeanCurvatureComparison}
  \uses{prop:pet-ch7-volume-form-laplacian-trace,def:pet-ch7-polar-density}
  If $(M,g)$ has $\Ric\ge(n-1)k$ for some $k\in\mathbb{R}$, then
  \[
    \Delta r\le(n-1)\frac{\snk_k'(r)}{\snk_k(r)},
    \qquad
    \partial_r\Big(\frac{\lambda}{\lambda_k}\Big)\le0,
    \qquad
    \lambda(r,\theta)\le\lambda_k(r)=\snk_k^{n-1}(r).
  \]
  \source{petersen:page-0293}
\end{lemma}
\begin{proof}
  \uses{lem:pet-ch7-ricci-comparison-mean-curvature}
  Divide (tr2) by $n-1$: with $\rho=\Delta r/(n-1)$,
  \[
    \partial_r\rho+\rho^2\le\frac{1}{n-1}\Big(\partial_r\Delta r+|\Hess r|^2\Big)
    =-\frac{\Ric(\partial_r,\partial_r)}{n-1}\le-k,
  \]
  using \cref{prop:pet-ch7-volume-form-laplacian-trace} and the hypothesis
  $\Ric\ge(n-1)k$. Thus $\rho$ solves the Riccati differential inequality
  $\dot\rho+\rho^2\le-k$ along $\partial_r$, with the same singular behavior
  $\rho(r)=1/r+O(r)$ as $r\to0$ as the comparison solution
  $\rho_k(t)=\snk_k'(t)/\snk_k(t)$, which solves $\dot\rho_k+\rho_k^2=-k$
  exactly. By the Riccati comparison estimate of Chapter 6 (Corollary 6.4.2),
  $\rho(r)\le\rho_k(r)$, i.e.
  $\Delta r\le(n-1)\snk_k'(r)/\snk_k(r)$.

  For the second inequality, \cref{def:pet-ch7-polar-density} gives
  $\partial_r\lambda=\lambda\Delta r\le(n-1)\dfrac{\snk_k'(r)}{\snk_k(r)}\lambda$
  (using $\lambda>0$ and the first inequality) and, exactly,
  $\partial_r\lambda_k=(n-1)\dfrac{\snk_k'(r)}{\snk_k(r)}\lambda_k$. Since
  $\lambda_k>0$ on the relevant range of $r$, the quotient rule gives
  \[
    \partial_r\Big(\frac{\lambda}{\lambda_k}\Big)
    =\frac{\lambda_k\,\partial_r\lambda-\lambda\,\partial_r\lambda_k}{\lambda_k^2}
    \le\frac{\lambda_k\cdot(n-1)\frac{\snk_k'}{\snk_k}\lambda
      -\lambda\cdot(n-1)\frac{\snk_k'}{\snk_k}\lambda_k}{\lambda_k^2}=0.
  \]

  Finally, near $r=0$ the polar density of any metric has the expansion
  $\lambda(r,\theta)=r^{n-1}+O(r^n)$ (the leading term of the volume of a
  small geodesic ball, from the Jacobi field expansion of $\exp_p$), and
  likewise $\snk_k(r)=r+O(r^3)$ gives
  $\lambda_k(r)=\snk_k^{n-1}(r)=r^{n-1}+O(r^n)$; hence
  $\lim_{r\to0}\lambda(r,\theta)/\lambda_k(r)=1$. Combined with
  $\partial_r(\lambda/\lambda_k)\le0$, the ratio $\lambda/\lambda_k$ is
  nonincreasing from its limiting value $1$ at $r=0$, so
  $\lambda(r,\theta)\le\lambda_k(r)$ for all $r$ in the domain.
\end{proof}

\begin{lemma}[Lemma 7.1.3]
  \label{lem:pet-ch7-absolute-volume-comparison}
  \lean{PetersenLib.absoluteVolumeComparison}
  \uses{lem:pet-ch7-ricci-comparison-mean-curvature}
  If $(M,g)$ has $\Ric\ge(n-1)k$, then $\vol B(p,r)\le v(n,k,r)$, where
  $v(n,k,r)$ is the volume of a metric ball of radius $r$ in the constant
  curvature space form $S^n_k$.
  \source{petersen:page-0294}
\end{lemma}
\begin{proof}
  \uses{lem:pet-ch7-absolute-volume-comparison}
  In exponential polar coordinates at $p$,
  \[
    \vol B(p,r)
    =\int_{S^{n-1}\cap B(0,r)}\lambda(r,\theta)\,dr\wedge\vol_{n-1}
    \le\int_{S^{n-1}\cap B(0,r)}\lambda_k(r)\,dr\wedge\vol_{n-1}
    \le\int_{B(0,r)}\vol_k
    =v(n,k,r),
  \]
  the first inequality by
  \cref{lem:pet-ch7-ricci-comparison-mean-curvature} and the second by
  extending the domain of integration from (a subset of) the segment domain
  at $p$ to all of the model ball $B(0,r)$ in $S^n_k$, using that the
  integrand $\lambda_k(r)$ is nonnegative.
\end{proof}

With a little more work, the absolute bound of
\cref{lem:pet-ch7-absolute-volume-comparison} improves to a relative
statement — the \emph{relative volume comparison} estimate, used repeatedly
throughout the rest of the book.

\begin{lemma}[Lemma 7.1.4 — Relative Volume Comparison, Bishop, 1964, and Gromov, 1980]
  \label{lem:pet-ch7-relative-volume-comparison}
  \lean{PetersenLib.relativeVolumeComparison}
  \uses{lem:pet-ch7-ricci-comparison-mean-curvature}
  Let $(M,g)$ be complete with $\Ric\ge(n-1)k$. The volume ratio
  \[
    r\mapsto\frac{\vol B(p,r)}{v(n,k,r)}
  \]
  is a nonincreasing function of $r$, with limit $1$ as $r\to0$.
  \source{petersen:page-0294,petersen:page-0295}
\end{lemma}
\begin{proof}
  \uses{lem:pet-ch7-relative-volume-comparison}
  Work in exponential polar coordinates at $p$. The polar density $\lambda$
  of $(M,g)$ is defined a priori only on the (star-shaped) segment domain in
  $T_pM\cong\mathbb{R}^n$; extend it by $\lambda=0$ outside. The comparison
  density $\lambda_k$ is defined on all of $\mathbb{R}^n$ when $k\le0$ and on
  $B(0,\pi/\sqrt k)$ when $k>0$; in the latter case extend $\lambda_k=0$
  outside $B(0,\pi/\sqrt k)$, and note that Myers' diameter theorem (Theorem
  6.3.3) forces $\lambda=0$ outside $B(0,\pi/\sqrt k)$ as well when $k>0$
  (since $M$ then has diameter $\le\pi/\sqrt k$, so no point of the segment
  domain lies at distance $>\pi/\sqrt k$ from $p$); it therefore suffices to
  consider $r<\pi/\sqrt k$ in that case. In either case,
  $0\le\lambda(r,\theta)\le\lambda_k(r)$ for all $(r,\theta)$, by
  \cref{lem:pet-ch7-ricci-comparison-mean-curvature}.

  Write $N(R)=\int_0^R\int_{S^{n-1}}\lambda\,dr\wedge\vol_{n-1}=\vol B(p,R)$
  and $D(R)=\int_0^R\int_{S^{n-1}}\lambda_k\,dr\wedge\vol_{n-1}=v(n,k,R)$.
  Differentiating the quotient,
  \[
    \frac{d}{dR}\frac{N(R)}{D(R)}=\frac{N'(R)D(R)-N(R)D'(R)}{D(R)^2},
  \]
  where $N'(R)=\int_{S^{n-1}}\lambda(R,\theta)\,\vol_{n-1}$ and
  $D'(R)=\int_{S^{n-1}}\lambda_k(R)\,\vol_{n-1}$. Writing
  $D(R)=\int_0^RD'(r)\,dr$, $N(R)=\int_0^RN'(r)\,dr$, this becomes
  \[
    \frac{d}{dR}\frac{N(R)}{D(R)}
    =D(R)^{-2}\int_0^R\Big(N'(R)D'(r)-D'(R)N'(r)\Big)\,dr,
  \]
  so it suffices that $N'(R)D'(r)\le D'(R)N'(r)$ for all $r\le R$, i.e.\
  (since $D'(r)=\omega_{n-1}\lambda_k(r)>0$) that
  \[
    r\mapsto\frac{N'(r)}{D'(r)}
    =\frac{1}{\omega_{n-1}}\int_{S^{n-1}}\frac{\lambda(r,\theta)}{\lambda_k(r)}\,\vol_{n-1}
  \]
  is nonincreasing. By \cref{lem:pet-ch7-ricci-comparison-mean-curvature},
  $\partial_r(\lambda(r,\theta)/\lambda_k(r))\le0$ pointwise in $\theta$, so
  the average over $\theta\in S^{n-1}$ is nonincreasing in $r$ as well. Hence
  $N(R)/D(R)=\vol B(p,R)/v(n,k,R)$ is nonincreasing in $R$.

  Finally, $\lim_{r\to0}\lambda(r,\theta)/\lambda_k(r)=1$ uniformly in
  $\theta$ (established in the proof of
  \cref{lem:pet-ch7-ricci-comparison-mean-curvature}), so
  $N'(r)/D'(r)\to1$ as $r\to0$; since also $N(r),D(r)\to0$ with $D(r)>0$ for
  $r>0$, this gives $N(r)/D(r)\to1$ as $r\to0$ as well.
\end{proof}

\subsection{The Maximum Principle}

We explain how one can assign second derivatives to functions at points where
the function is not smooth, using \emph{support functions} (also called
barrier functions in PDE theory) rather than the generalized gradients of
Lipschitz functions.

\begin{lemma}[Lemma 7.1.5]
  \label{lem:pet-ch7-hessian-touching-comparison}
  \lean{PetersenLib.hessianTouchingComparison}
  If $f,h:(M,g)\to\mathbb{R}$ are $C^2$ functions such that $f(p)=h(p)$ and
  $f(x)\ge h(x)$ for all $x$ near $p$, then
  \[
    \nabla f(p)=\nabla h(p),\qquad \Hess f|_p\ge \Hess h|_p,\qquad \Delta f(p)\ge \Delta h(p).
  \]
  \source{petersen:page-0295,petersen:page-0296}
\end{lemma}
\begin{proof}
  \uses{lem:pet-ch7-hessian-touching-comparison}
  If $(M,g)\subset(\mathbb{R},g_\mathbb{R})$ this is standard single-variable calculus.
  In general let $c:(-\varepsilon,\varepsilon)\to M$ be a curve with $c(0)=p$. Applying
  the one-dimensional statement to $f\circ c,\,h\circ c$ (which agree at $0$ and
  satisfy $f\circ c\ge h\circ c$ near $0$) gives
  \[
    df(\dot c(0))=dh(\dot c(0)),\qquad \Hess f(\dot c(0),\dot c(0))\ge \Hess h(\dot c(0),\dot c(0)).
  \]
  Letting $v=\dot c(0)$ range over all of $T_pM$ (every $v$ is $\dot c(0)$ for
  some geodesic $c$ with $c(0)=p$) yields $\nabla f(p)=\nabla h(p)$ and
  $\Hess f|_p\ge\Hess h|_p$; taking traces gives $\Delta f(p)\ge \Delta h(p)$.
\end{proof}

\begin{definition}[support functions; barrier-sense Hessian and Laplacian bounds]
  \label{def:pet-ch7-support-function}
  \lean{PetersenLib.SupportFunctionBelow,PetersenLib.barrierHessBound}
  \uses{lem:pet-ch7-hessian-touching-comparison}
  By \cref{lem:pet-ch7-hessian-touching-comparison}, a $C^2$ function
  $f:M\to\mathbb{R}$ has $\Hess f|_p\ge B$ for a symmetric bilinear form $B$ on $T_pM$
  (respectively $\Delta f(p)\ge a\in\mathbb{R}$) if and only if for every
  $\varepsilon>0$ there is a function $f_\varepsilon$ defined near $p$, called a
  \emph{support function from below} for $f$ at $p$, with
  \begin{enumerate}
    \item[(1)] $f_\varepsilon(p)=f(p)$;
    \item[(2)] $f(x)\ge f_\varepsilon(x)$ near $p$;
    \item[(3)] $\Hess f_\varepsilon|_p\ge B-\varepsilon g|_p$ (respectively
      $\Delta f_\varepsilon(p)\ge a-\varepsilon$).
  \end{enumerate}
  \emph{Support functions from above} are defined symmetrically and give upper
  bounds for $\Hess f$ and $\Delta f$. For a merely continuous $f:(M,g)\to\mathbb{R}$
  one \emph{defines} $\Hess f|_p\ge B$ (respectively $\Delta f(p)\ge a$), said to
  hold in the \emph{support} or \emph{barrier sense}, to mean that such smooth
  support functions $f_\varepsilon$ from below exist for every $\varepsilon>0$.
  A continuous $f:(M,g)\to\mathbb{R}$ is \emph{convex} if $\Hess f\ge0$ everywhere in
  this sense; this matches the classical definition when $(M,g)\subset(\mathbb{R},g_\mathbb{R})$.
  \source{petersen:page-0296}
\end{definition}

\begin{theorem}[Theorem 7.1.6]
  \label{thm:pet-ch7-convex-constant-near-max}
  \lean{PetersenLib.convexHessNonnegConstantNearMax}
  \uses{def:pet-ch7-support-function}
  If $f:(M,g)\to\mathbb{R}$ is continuous with $\Hess f\ge0$ everywhere, then $f$ is
  constant near any local maximum. In particular, $f$ cannot have a global
  maximum unless $f$ is constant.
  \source{petersen:page-0296}
\end{theorem}

\begin{definition}[subharmonic and superharmonic functions]
  \label{def:pet-ch7-subharmonic-superharmonic}
  \lean{PetersenLib.IsSubharmonic,PetersenLib.IsSuperharmonic}
  \uses{def:pet-ch7-support-function}
  A continuous $f:(M,g)\to\mathbb{R}$ with $\Delta f\ge0$ everywhere (in the support
  sense) is \emph{subharmonic}; if $\Delta f\le0$ it is \emph{superharmonic}.
  \source{petersen:page-0296}
\end{definition}

\begin{theorem}[Theorem 7.1.7 — The Strong Maximum Principle]
  \label{thm:pet-ch7-strong-maximum-principle}
  \lean{PetersenLib.strongMaximumPrinciple}
  \uses{def:pet-ch7-subharmonic-superharmonic,thm:pet-ch7-convex-constant-near-max}
  If $f:(M,g)\to\mathbb{R}$ is continuous and subharmonic, then $f$ is constant in a
  neighborhood of every local maximum. In particular, if $f$ has a global
  maximum then $f$ is constant.
  \source{petersen:page-0296,petersen:page-0297,petersen:page-0298}
\end{theorem}
\begin{proof}
  \uses{thm:pet-ch7-strong-maximum-principle,def:pet-ch7-support-function,lem:pet-ch7-hessian-touching-comparison}
  First suppose $\Delta f>0$ everywhere. If $f$ had a local maximum at
  $p$, there would be a smooth support function $f_\varepsilon$ from below with
  $f_\varepsilon(p)=f(p)$, $f_\varepsilon\le f$ near $p$, and
  $\Delta f_\varepsilon(p)>0$. Conditions (1)-(2) force $f_\varepsilon$ also to
  have a local maximum at $p$, so $\Hess f_\varepsilon(p)\le0$ and hence
  $\Delta f_\varepsilon(p)\le0$, a contradiction. Thus $f$ has no local maxima
  at all in this case.

  Now assume only $\Delta f\ge0$, and let $p\in M$ be a local maximum. For
  $r<\operatorname{inj}(p)$ small enough, $f|_{B(p,r)}$ has a global maximum at
  $p$. If $f$ is constant on $B(p,r)$ we are done; otherwise, after possibly
  shrinking $r$, there is $x_0\in\partial B(p,r)$ with $f(x_0)\ne f(p)$. Put
  \[
    V=\{x\in\partial B(p,r)\mid f(x)=f(p)\}.
  \]
  Choose an open disc $U\subset\partial B(p,r)$ containing $V$ and a function
  $\phi$ with $\phi(p)=0$, $\phi<0$ on $U$, and $\nabla\phi\ne0$ on
  $\overline{B}(p,r)$: take $\phi=x^1$ in a coordinate system $(x^1,\dots,x^n)$
  centered at $p$ chosen so that $U$ lies in $\{x^1<0\}$. For $\alpha$ large
  enough,
  \[
    h=e^{\alpha\phi}-1
  \]
  satisfies $\Delta h=\alpha e^{\alpha\phi}(\alpha|\nabla\phi|^2+\Delta\phi)>0$
  on $\overline{B}(p,r)$, while $h(p)=0$ and $h<0$ on $V\subset U$ (since
  $\phi<0$ there).

  Consider $\bar f=f+\delta h$ on $\overline{B}(p,r)$. For $\delta>0$ small
  enough, $h<0$ on $V$ and continuity of $f-f(p)$ (which vanishes only on $V$
  among boundary maximizers) force
  \[
    \bar f(p)=f(p)>\max\{\bar f(x)\mid x\in\partial B(p,r)\},
  \]
  so $\bar f$ has a local maximum at some interior point $q\in B(p,r)$. Fix
  $\varepsilon>0$ and let $f_\varepsilon$ be a support function from below for
  $f$ at $q$; then $f_\varepsilon+\delta h$ is a support function from below
  for $\bar f$ at $q$, and
  \[
    \Delta(f_\varepsilon+\delta h)(q)\ge-\varepsilon+\delta\,\Delta h(q),
  \]
  which is positive for $\varepsilon$ small since $\Delta h(q)>0$. This
  contradicts the first paragraph applied to $\bar f$ (which has $\Delta\bar
  f>0$ in the support sense near its interior maximum $q$), so no such $x_0$
  exists and $f$ is constant on $B(p,r)$.
\end{proof}

\begin{definition}[linear and harmonic functions]
  \label{def:pet-ch7-linear-harmonic-function}
  \lean{PetersenLib.IsLinearFunction,PetersenLib.IsHarmonicFunction}
  \uses{def:pet-ch7-support-function}
  A continuous $f:(M,g)\to\mathbb{R}$ is \emph{linear} if $\Hess f\equiv0$, i.e.\
  both $\Hess f\ge0$ and $\Hess f\le0$ hold everywhere in the support sense.
  Then $f\circ c$ is both convex and concave along every geodesic $c$, so
  $(f\circ c)(t)=f(c(0))+at$; equivalently
  $f\circ\exp_p(x)=f(p)+g(v_p,x)$ for a fixed $v_p\in T_pM$, and $f$ is
  automatically $C^\infty$ with $\nabla f|_p=v_p$. More generally, $f$ is
  \emph{harmonic} if it is continuous with $\Delta f=0$ (in the support sense).
  \source{petersen:page-0298}
\end{definition}

\begin{theorem}[Theorem 7.1.8 — Regularity of harmonic functions]
  \label{thm:pet-ch7-harmonic-regularity}
  \lean{PetersenLib.harmonicFunctionRegularity}
  \uses{def:pet-ch7-linear-harmonic-function}
  If $f:(M,g)\to\mathbb{R}$ is continuous and harmonic in the weak (support) sense,
  then $f$ is smooth.
  \source{petersen:page-0298,petersen:page-0299}
\end{theorem}
\begin{proof}
  \uses{thm:pet-ch7-harmonic-regularity,thm:pet-ch7-strong-maximum-principle}
  Fix $p\in M$ and a neighborhood $\Omega$ of $p$ with smooth boundary,
  contained in a coordinate chart. By standard, nontrivial elliptic PDE theory
  the Dirichlet problem
  \[
    \Delta u=0,\qquad u|_{\partial\Omega}=f|_{\partial\Omega}
  \]
  has a solution $u$ that is smooth on the interior of $\Omega$. Consider
  $u-f$ and $f-u$ on $\Omega$: both vanish on $\partial\Omega$, and each is
  subharmonic (being the difference of a harmonic and a harmonic-in-the-weak-sense
  function, so of two functions with $\Delta\ge0$, hence itself has $\Delta\ge0$
  in the support sense). If both are $\le0$ on $\Omega$, they both vanish and
  $f=u$ is smooth near $p$. Otherwise one of them is positive somewhere in
  $\Omega$; being subharmonic and vanishing on $\partial\Omega$, it then attains
  an interior global maximum, so by \cref{thm:pet-ch7-strong-maximum-principle}
  it is constant on $\Omega$ — but a nonnegative constant vanishing on the
  boundary must be $0$, a contradiction. Hence $f=u$ is smooth near $p$, and
  since $p$ was arbitrary, $f$ is smooth on $M$.
\end{proof}

\subsection{Geometric Laplacian Comparison}

The idea of using support functions to estimate the Laplacian is particularly
convenient for geometric applications, since distance functions always have
support functions from above (obtained from the triangle inequality along a
fixed segment to the base point).

\begin{lemma}[Lemma 7.1.9 — Calabi, 1958]
  \label{lem:pet-ch7-laplacian-comparison-calabi}
  \lean{PetersenLib.laplacianComparisonCalabi}
  \uses{def:pet-ch7-support-function}
  If $(M,g)$ is complete and $\Ric(M,g)\ge(n-1)k$, then any distance function
  $r(x)=|xp|$ satisfies
  \[
    \Delta r(x)\le(n-1)\frac{\snk_k'(r(x))}{\snk_k(r(x))}.
  \]
  \source{petersen:page-0298,petersen:page-0299}
\end{lemma}
\begin{proof}
  \uses{lem:pet-ch7-laplacian-comparison-calabi,lem:pet-ch7-ricci-comparison-mean-curvature}
  By \cref{lem:pet-ch7-ricci-comparison-mean-curvature} (mean curvature
  comparison), the stated inequality already holds pointwise wherever $r$ is
  smooth. In general fix $q\in M$ and a unit-speed segment $\sigma:[0,L]\to M$ with
  $\sigma(0)=p$, $\sigma(L)=q$. By the triangle inequality,
  \[
    r_\varepsilon(x)=\varepsilon+|\sigma(\varepsilon)x|
  \]
  is a support function from above for $r$ at $q$: $r_\varepsilon(q)=r(q)$ and
  $r_\varepsilon(x)\ge r(x)$ near $q$ (a segment from $\sigma(\varepsilon)$ to
  $x$ together with the segment from $p$ to $\sigma(\varepsilon)$ need not be
  the shortest path from $p$ to $x$).

  Suppose first each $r_\varepsilon$ is smooth at $q$. Applying the smooth case
  to the distance function based at $\sigma(\varepsilon)$,
  \[
    \Delta r_\varepsilon(q)\le(n-1)\frac{\snk_k'(r_\varepsilon(q))}{\snk_k(r_\varepsilon(q))}
    =(n-1)\frac{\snk_k'(r(q)-\varepsilon)}{\snk_k(r(q)-\varepsilon)}
    \searrow(n-1)\frac{\snk_k'(r(q))}{\snk_k(r(q))}
  \]
  as $\varepsilon\to0$, since $t\mapsto \snk_k'(t)/\snk_k(t)$ is decreasing.
  Since $\Delta r(q)\le\Delta r_\varepsilon(q)$ for every $\varepsilon$ (as
  $r_\varepsilon$ is a support function from above), letting $\varepsilon\to0$
  gives the claimed bound at $q$.

  It remains to justify smoothness of $r_\varepsilon$ at $q$ for $\varepsilon>0$
  small. If $r_\varepsilon$ fails to be smooth at $q$, then (by the criterion
  for smoothness of distance functions) either (1) there are two distinct
  segments from $\sigma(\varepsilon)$ to $q$, or (2) $q$ is a critical value of
  $\exp_{\sigma(\varepsilon)}$. Case (1) would produce a broken curve of length
  $L$ from $p$ to $q$ (concatenating $\sigma|_{[0,\varepsilon]}$ with either of
  the two segments), forced by minimality of $L=|pq|$ to itself be a
  (necessarily smooth, hence unbroken) segment; this rules out two distinct
  segments from $\sigma(\varepsilon)$ to $q$ unless they coincide, contradicting
  (1). So case (2) must hold: $q$ is a critical value of
  $\exp_{\sigma(\varepsilon)}$, which in turn forces $q$ to be a critical value
  of $\exp_p$ as well (the segment $\sigma$ concatenated with a Jacobi field
  vanishing at $\sigma(\varepsilon)$ produces a Jacobi field along $\sigma$
  vanishing at $p$ that also degenerates at $q$), contradicting that $r=|{\cdot}p|$
  is smooth at points that are not critical values of $\exp_p$ and that
  $q\ne p$ lies on the minimizing segment $\sigma$, a locus on which $\exp_p$ is
  nonsingular up to $q$. This contradiction shows $r_\varepsilon$ is smooth at
  $q$ for $\varepsilon$ small, completing the proof.
\end{proof}

\subsection{The Segment, Poincaré, and Sobolev Inequalities}

We use the volume comparison results of §7.1.2 to prove analytic inequalities
used later in the book (Chapter 9). Throughout, for a bounded domain $B\subset
M$ and $u:M\to\mathbb{R}$, write
\[
  \|u\|_{p,B}=\left(\frac{1}{\vol B}\int_B|u|^p\,\vol\right)^{1/p},\qquad
  u_B=\frac{1}{\vol B}\int_Bu\,\vol.
\]

\begin{theorem}[Theorem 7.1.10 — The Segment Inequality, Cheeger and Colding, 1996]
  \label{thm:pet-ch7-segment-inequality}
  \lean{PetersenLib.segmentInequality}
  Assume $(M,g)$ has $\Ric\ge(n-1)k$, $k\le0$. Let $f:M\to[0,\infty)$ and
  $A,B\subset W\subset M$ with $\diam W\le D$. For each $x\in A$, $y\in B$ fix a
  segment $c_{x,y}:[0,1]\to M$ from $x$ to $y$, and suppose $c_{x,y}(t)\in W$
  for all $x\in A$, $y\in B$, $t\in[0,1]$. Then
  \[
    \int_A\int_B\int_0^1 f\circ c_{x,y}(t)\,dt\,\vol_x\,\vol_y
    \le C(\vol A+\vol B)\int_Wf\,\vol,
  \]
  where $C=C(n,kD^2)$.
  \source{petersen:page-0300}
\end{theorem}
\begin{proof}
  \uses{thm:pet-ch7-segment-inequality,lem:pet-ch7-ricci-comparison-mean-curvature}
  Let
  \[
    C=\max_{R\le D}\frac{\snk_k^{n-1}(R)}{\snk_k^{n-1}(R/2)},
  \]
  which is $2^{n-1}$ when $k=0$ and
  $\sinh^{n-1}(\sqrt{-k}D)/\sinh^{n-1}(\tfrac12\sqrt{-k}D)$ otherwise; in either
  case $C=C(n,kD^2)$.

  Fix $x\in A$ and use polar coordinates centered at $x$ with polar density
  $\lambda$. For $t\ge\tfrac12$ the map $y\mapsto c_{x,y}(t)$ rescales radial
  distance from $x$ by the factor $t$ along each segment, so writing
  $y=(r,\theta)$ in these polar coordinates,
  \[
    \int_Bf\circ c_{x,y}(t)\,\vol_y
    =\int_Bf\circ c_{x,y}(t)\,\lambda(y)\,dr\,\vol_{n-1}
    =\int_Bf\circ c_{x,y}(t)\,\lambda(c_{x,y}(t))\,
      \frac{\lambda(y)}{\lambda(c_{x,y}(t))}\,dr\,\vol_{n-1}.
  \]
  Mean curvature comparison (\cref{lem:pet-ch7-ricci-comparison-mean-curvature})
  shows $r\mapsto\lambda(r,\theta)$ grows no
  faster than the comparison density, so along the ray from $x$ through $y$ at
  parameter $t\ge\tfrac12$ the ratio $\lambda(y)/\lambda(c_{x,y}(t))$ is bounded
  by the comparison ratio $\snk_k^{n-1}(r)/\snk_k^{n-1}(tr)\le
  \snk_k^{n-1}(D)/\snk_k^{n-1}(D/2)\le C$. Hence
  \[
    \int_Bf\circ c_{x,y}(t)\,\vol_y
    \le C\int_Bf\circ c_{x,y}(t)\,\lambda(c_{x,y}(t))\,dr\,\vol_{n-1}
    \le C\int_Wf\,\vol,
  \]
  the last step being a change of variables to the point $c_{x,y}(t)\in W$.
  Integrating over $t\in[\tfrac12,1]$ and $x\in A$,
  \[
    \int_A\int_B\int_{1/2}^1f\circ c_{x,y}(t)\,dt\,\vol_x\,\vol_y
    \le\tfrac12C\,\vol A\int_Wf\,\vol.
  \]
  By the symmetric argument (swapping the roles of $A,B$ and using
  $t\le\tfrac12$, so $1-t\ge\tfrac12$ rescales from $y$ instead of $x$),
  \[
    \int_B\int_A\int_0^{1/2}f\circ c_{x,y}(t)\,dt\,\vol_x\,\vol_y
    \le\tfrac12C\,\vol B\int_Wf\,\vol.
  \]
  Adding the two estimates gives the segment inequality.
\end{proof}

\begin{corollary}[Corollary 7.1.11]
  \label{cor:pet-ch7-weak-poincare}
  \lean{PetersenLib.weakPoincareInequality}
  \uses{thm:pet-ch7-segment-inequality}
  Assume $(M,g)$ has $\Ric\ge(n-1)k$, $k\le0$. Any smooth $u:M\to[0,\infty)$
  satisfies
  \[
    \|u-u_{B(p,R)}\|_{1,B(p,R)}\le4C^2R\,\|du\|_{1,B(p,2R)},
  \]
  where $R\le D$ and $C=C(n,kD^2)$ is as in \cref{thm:pet-ch7-segment-inequality}.
  \source{petersen:page-0301,petersen:page-0302}
\end{corollary}
\begin{proof}
  \uses{cor:pet-ch7-weak-poincare,thm:pet-ch7-segment-inequality,lem:pet-ch7-relative-volume-comparison}
  Write $B=B(p,R)$, $W=B(p,2R)$. For $x,y\in B$ let $c_{x,y}$ be a fixed
  segment from $x$ to $y$; since $\diam B\le2R$, $c_{x,y}$ lies in $W$.
  Then
  \begin{align*}
    \int_B|u-u_B|
    &=\int_B\left|\frac{1}{\vol B}\int_B(u(x)-u(y))\,\vol_y\right|\vol_x\\
    &\le\frac{1}{\vol B}\int_B\int_B|u(x)-u(y)|\,\vol_y\,\vol_x\\
    &\le\frac{1}{\vol B}\int_B\int_B\int_0^1|xy|\,|du|(c_{x,y}(t))\,dt\,\vol_y\,\vol_x\\
    &\le\frac{2R}{\vol B}\int_B\int_B\int_0^1|du|(c_{x,y}(t))\,dt\,\vol_y\,\vol_x,
  \end{align*}
  using $|u(x)-u(y)|\le\int_0^1|du|(c_{x,y}(t))\,|\dot c_{x,y}(t)|\,dt$ and
  $|\dot c_{x,y}|\equiv|xy|$, then $|xy|\le\diam B\le2R$. Applying
  \cref{thm:pet-ch7-segment-inequality} with $A=B=B(p,R)$, $W=B(p,2R)$,
  $f=|du|$ and $D=2R$,
  \[
    \int_B\int_B\int_0^1|du|(c_{x,y}(t))\,dt\,\vol_x\,\vol_y
    \le2C\,\vol B\int_W|du|\,\vol,
  \]
  so
  \[
    \|u-u_B\|_{1,B}
    \le\frac{2R}{\vol B}\cdot2C\,\vol B\int_W|du|\,\vol
    =4CR\,\frac{\vol W}{\vol B}\,\|du\|_{1,W}.
  \]
  Relative volume comparison (\cref{lem:pet-ch7-relative-volume-comparison})
  bounds $\vol W/\vol B=\vol B(p,2R)/\vol
  B(p,R)\le v(n,k,2D)/v(n,k,D)$, a constant depending only on $n,kD^2$, which
  may be absorbed into the constant $C$ (replacing $4C$ by $4C^2$) to give the
  stated inequality.
\end{proof}

\begin{remark}[Remark 7.1.12]
  \label{rem:pet-ch7-upper-gradient}
  \lean{PetersenLib.upperGradientPoincare}
  \uses{cor:pet-ch7-weak-poincare}
  \Cref{cor:pet-ch7-weak-poincare} holds for any measurable $u$ with a function
  $G$ in place of $|du|$, provided $|u(x)-u(y)|\le\int_0^1G(c(t))\,|\dot
  c(t)|\,dt$ for all $c\in\Omega_{x,y}$ (paths from $x$ to $y$); such $G$ is an
  \emph{upper gradient} for $u$.
  \source{petersen:page-0302}
\end{remark}

\begin{definition}[maximal function]
  \label{def:pet-ch7-maximal-function}
  \lean{PetersenLib.maximalFunction}
  For $u:M\to\mathbb{R}$ and $D>0$, the \emph{maximal function} of $u$ (relative to
  the scale $D$) is
  \[
    M(u)(x)=\sup_{R\in(0,D]}\frac{1}{\vol B(x,R)}\int_{B(x,R)}|u|\,\vol.
  \]
  \source{petersen:page-0302}
\end{definition}

\begin{theorem}[Theorem 7.1.14 — Maximal Function Theorem]
  \label{thm:pet-ch7-maximal-function-theorem}
  \lean{PetersenLib.maximalFunctionTheorem}
  \uses{def:pet-ch7-maximal-function}
  Assume $(M,g)$ has $\Ric\ge(n-1)k$. There is $C=C(n,kD^2)$ such that for all
  $t>0$,
  \[
    t\,\vol\{M(u)>t\}\le C\int|u|\,\vol.
  \]
  \source{petersen:page-0303}
\end{theorem}
\begin{proof}
  \uses{thm:pet-ch7-maximal-function-theorem,lem:pet-ch7-relative-volume-comparison,rem:pet-ch7-ex-5}
  For each $x\in\{M(u)>t\}$ there is $R_x\le D$ with
  \[
    t\,\vol B(x,R_x)<\int_{B(x,R_x)}|u|\,\vol.
  \]
  By the basic covering lemma (\cref{rem:pet-ch7-ex-5}: given a bounded
  positive function $R:X\to(0,D]$ on a separable metric space there is a
  countable $A\subset X$ with the balls $B(p,R(p))$, $p\in A$, pairwise
  disjoint and $X=\bigcup_{p\in A}B(p,5R(p))$), we may cover $\{M(u)>t\}$ by
  balls $B(x_i,5R_{x_i})$ with the $B(x_i,R_{x_i})$ pairwise disjoint.
  Relative volume comparison (\cref{lem:pet-ch7-relative-volume-comparison})
  applied to the two radii $R\le5R$ gives
  \[
    \frac{\vol B(x,5R)}{\vol B(x,R)}\le\frac{v(n,k,5D)}{v(n,k,D)}=:C=C(n,kD^2)
  \]
  for all $x$ and $R\le D$. Hence
  \begin{align*}
    t\,\vol\{M(u)>t\}
    &\le\sum_it\,\vol B(x_i,5R_{x_i})
    \le C\sum_it\,\vol B(x_i,R_{x_i})\\
    &<C\sum_i\int_{B(x_i,R_{x_i})}|u|\,\vol
    \le C\int|u|\,\vol,
  \end{align*}
  the last step using pairwise disjointness of the $B(x_i,R_{x_i})$.
\end{proof}

\begin{theorem}[Theorem 7.1.15]
  \label{thm:pet-ch7-weak-poincare-sobolev-estimate}
  \lean{PetersenLib.weakPoincareSobolevEstimate}
  \uses{cor:pet-ch7-weak-poincare,thm:pet-ch7-maximal-function-theorem}
  Assume $(M,g)$ has $\Ric\ge(n-1)k$, $k\le0$, and $\diam M\le D$. For smooth
  $u:M\to[0,\infty)$,
  \[
    t^{\frac{n}{n-1}}\,\vol\{u-u_{B(x,R)}>t\}
    \le C R^{\frac{n}{n-1}}\,\vol M\,\|du\|_1^{\frac{n}{n-1}},
  \]
  where $C=C(n,kD^2)$.
  \source{petersen:page-0303,petersen:page-0304,petersen:page-0305}
\end{theorem}
\begin{proof}
  \uses{thm:pet-ch7-weak-poincare-sobolev-estimate,cor:pet-ch7-weak-poincare,thm:pet-ch7-maximal-function-theorem,lem:pet-ch7-relative-volume-comparison}
  We prove the case $R=D$; the general case follows by the same argument
  restricted to $B(x,R)$. Fix $x\in M$, set $R_i=2^{-i}D$ and $B_i=B(x,R_i)$,
  so $M=B_0$. Continuity of $u$ gives $u(x)=\lim_iu_{B_i}$, so
  \[
    |u(x)-u_{B_0}|
    \le\sum_{i=0}^\infty|u_{B_i}-u_{B_{i+1}}|
    \le\sum_{i=0}^\infty\|u-u_{B_i}\|_{1,B_{i+1}}
    \le\sum_{i=0}^\infty\frac{\vol B_i}{\vol B_{i+1}}\|u-u_{B_i}\|_{1,B_i}.
  \]
  Relative volume comparison (\cref{lem:pet-ch7-relative-volume-comparison})
  bounds $\vol B_i/\vol B_{i+1}\le C$, and
  \cref{cor:pet-ch7-weak-poincare} bounds $\|u-u_{B_i}\|_{1,B_i}\le
  4C^2R_i\|du\|_{1,B_{i-1}}$, so
  \[
    |u(x)-u_{B_0}|\le2C^3\sum_{i=0}^\infty R_i\,\|du\|_{1,B_{i-1}}.
  \]
  It thus suffices to bound $\vol\{\sum_iR_i\|du\|_{1,B_{i-1}}>t\}$. For
  $r>0$ split the sum at $R_i\le r$ versus $R_i>r$. The tail with $R_i\le r$ is
  controlled by the maximal function:
  \[
    \sum_{R_i\le r}R_i\|du\|_{1,B_{i-1}}
    \le\Big(\sum_{R_i\le r}R_i\Big)M(|du|)(x)\le2r\,M(|du|)(x).
  \]
  The head with $R_i>r$ is controlled crudely by $\vol M/\vol B_{i-1}\le
  C(D/R_{i-1})^n$ (relative volume comparison,
  \cref{lem:pet-ch7-relative-volume-comparison}), giving a geometric series in
  $i$ dominated by its first term, of size $\le C\,r^{1-n}D^n\|du\|_1$ up to a
  constant depending only on $n,kD^2$. Choosing
  $r=D\big(\|du\|_1/M(|du|)(x)\big)^{1/n}$ balances the two contributions and
  yields
  \[
    \sum_{i=0}^\infty R_i\|du\|_{1,B_{i-1}}
    \le CD\,\big(M(|du|)(x)\big)^{\frac{n-1}{n}}\,\|du\|_1^{1/n}.
  \]
  Consequently
  \[
    \vol\Big\{\sum_iR_i\|du\|_{1,B_{i-1}}>t\Big\}
    \le\vol\Big\{CD^{\frac{n}{n-1}}\|du\|_1^{\frac{1}{n-1}}M(|du|)(x)>t^{\frac{n}{n-1}}\Big\}
    \le t^{-\frac{n}{n-1}}CD^{\frac{n}{n-1}}\|du\|_1^{\frac1{n-1}}\int_M|du|\,\vol,
  \]
  the last inequality by \cref{thm:pet-ch7-maximal-function-theorem} applied to
  $|du|$ with threshold $t^{n/(n-1)}\big(CD^{n/(n-1)}\|du\|_1^{1/(n-1)}\big)^{-1}$.
  Since $\int_M|du|\,\vol=\vol M\cdot\|du\|_1$, this is
  $t^{-n/(n-1)}CD^{n/(n-1)}\vol M\,\|du\|_1^{n/(n-1)}$, which combined with
  $|u(x)-u_{B_0}|\le2C^3\sum_iR_i\|du\|_{1,B_{i-1}}$ gives the theorem after
  absorbing constants.
\end{proof}

\begin{theorem}[Theorem 7.1.13 — The Poincaré--Sobolev Inequality]
  \label{thm:pet-ch7-poincare-sobolev}
  \lean{PetersenLib.poincareSobolevInequality}
  \uses{thm:pet-ch7-weak-poincare-sobolev-estimate}
  Assume $(M,g)$ has $\Ric\ge(n-1)k$, $k\le0$. For all smooth $u:M\to[0,\infty)$
  and $\nu\in[1,\tfrac{n}{n-1}]$,
  \[
    \|u-u_{B(x,R)}\|_{\nu,B(x,R)}\le C(n,kD^2)\,R\,\|du\|_{1,B(x,R)},
  \]
  where $R\le D$.
  \source{petersen:page-0302,petersen:page-0305,petersen:page-0306,petersen:page-0307}
\end{theorem}
\begin{proof}
  \uses{thm:pet-ch7-poincare-sobolev,thm:pet-ch7-weak-poincare-sobolev-estimate}
  It suffices to treat $\nu=\tfrac{n}{n-1}$ and $B(x,R)=M$ with $D$ an upper
  bound for $\diam M$ (the general case is the same argument applied inside
  $B(x,R)$, and smaller $\nu\le n/(n-1)$ follows from Hölder's inequality on
  the probability measure $\vol(\cdot)/\vol B(x,R)$). Throughout, $C$ denotes a
  constant depending only on $n,kD^2$, possibly changing from line to line.

  First, two elementary reductions. For any constant $c$,
  \[
    \|u-u_M\|_p\le|c-u_M|+\|u-c\|_p\le2\|u-c\|_p,
  \]
  using $|c-u_M|=|(u-c)_M|\le\|u-c\|_p$ (Jensen); hence it suffices to bound
  $\inf_c\|u-c\|_p$ for a well-chosen $c$.

  Choose $m$ with $\vol\{u\ge m\}\ge\tfrac12\vol M$ and $\vol\{u\le m\}\ge
  \tfrac12\vol M$ (a median of $u$), and split $u-m=v^+-v^-$ with
  $v^+=\max\{u-m,0\}$, $v^-=\max\{m-u,0\}$; both satisfy $\vol\{v^\pm=0\}\ge
  \tfrac12\vol M$, and $|dv^+|+|dv^-|\le|du|$ pointwise (with $|dv^\pm|:=0$
  where $v^\pm=0$). Since
  $\|u-m\|_{n/(n-1)}\le\|v^+\|_{n/(n-1)}+\|v^-\|_{n/(n-1)}$, it suffices to
  bound $\|v\|_{n/(n-1)}\le CD\|dv\|_1$ for $v=v^\pm$.

  We first record: for $c$ with $\vol\{v=0\}\ge\tfrac12\vol M$,
  \[
    \vol\{v>t\}\le2\,\vol\{|v-c|>t/2\}
  \]
  for every $t>0$ and every such $c$ (if $t/2\le c$, then $\{c-v>t/2\}\subset
  \{v=0\}$ has volume $\ge\tfrac12\vol M$, so
  $\vol\{|v-c|>t/2\}\ge\vol\{c-v>t/2\}\ge\tfrac12\vol M\ge\tfrac12\vol\{v>t\}$;
  if $t/2>c\ge0$, then $\{v>t\}\subset\{v-c>t/2\}\subset\{|v-c|>t/2\}$
  directly). Taking $c=0$ in particular gives $\vol\{v>t\}\le2\vol\{v>t/2\}$
  trivially, but the useful form below is with $c$ ranging freely.

  For $0<a<b$ define the truncation
  \[
    v_a^b(x)=\begin{cases}b-a,&v(x)\ge b,\\v(x)-a,&a<v(x)\le b,\\0,&v(x)\le a,\end{cases}
  \]
  which has upper gradient $|dv|\cdot\chi_{\{a<v\le b\}}$ and satisfies
  $\vol\{v_a^b=0\}\ge\vol\{v\le a\}\ge\tfrac12\vol M$ whenever $a\ge0$ is such
  that $\{v\le a\}$ still has volume $\ge\tfrac12\vol M$ (true here since
  $\vol\{v=0\}\ge\tfrac12\vol M$ and $a\ge0$). Applying
  \cref{thm:pet-ch7-weak-poincare-sobolev-estimate} to $v_a^b$ and the bound of
  the previous paragraph with $c=0$,
  \[
    t^{\frac{n}{n-1}}\vol\{v_a^b>t\}
    \le2\cdot t^{\frac{n}{n-1}}\vol\{|v_a^b|>t/2\}
    \le C D^{\frac{n}{n-1}}\vol M\,\big\||dv|\chi_{\{a<v\le b\}}\big\|_1^{\frac{n}{n-1}}.
  \]

  Now decompose the integral of $v^{n/(n-1)}$ dyadically: with $k$ ranging over
  integers,
  \[
    \int v^{\frac{n}{n-1}}\,\vol
    \le\sum_k2^{\frac{kn}{n-1}}\vol\{2^{k-1}<v\le2^k\}
    \le\sum_k2^{\frac{kn}{n-1}}\vol\{v>2^{k-1}\}.
  \]
  Since $\{v>2^{k-1}\}\subset\{v_{2^{k-2}}^{2^{k-1}}>2^{k-2}\}$, taking
  $t=2^{k-2}$ in the truncated estimate above gives
  \[
    2^{\frac{kn}{n-1}}\vol\{v>2^{k-1}\}
    \le C\,2^{\frac{kn}{n-1}}2^{-\frac{kn}{n-1}}D^{\frac{n}{n-1}}\vol M\,
      \big\||dv|\chi_{\{2^{k-2}<v\le2^{k-1}\}}\big\|_1^{\frac{n}{n-1}}
    =CD^{\frac{n}{n-1}}\vol M\,\big\||dv|\chi_{\{2^{k-2}<v\le2^{k-1}\}}\big\|_1^{\frac{n}{n-1}}
  \]
  up to the absolute constant from the truncated estimate. Since the level
  sets $\{2^{k-2}<v\le2^{k-1}\}$ are pairwise disjoint in $k$, and
  $p\mapsto\ell^p$ norms decrease in $p$, summing over $k$ and using
  $\big(\sum_ka_k\big)^{n/(n-1)}\ge\sum_ka_k^{n/(n-1)}$ for $a_k\ge0$ in reverse
  (i.e. $\sum_ka_k^{n/(n-1)}\le\big(\sum_ka_k\big)^{n/(n-1)}$, which holds since
  $n/(n-1)\ge1$) with $a_k=\||dv|\chi_{\{2^{k-2}<v\le2^{k-1}\}}\|_1$ gives
  \[
    \sum_kCD^{\frac{n}{n-1}}\vol M\,a_k^{\frac{n}{n-1}}
    \le CD^{\frac{n}{n-1}}\vol M\Big(\sum_ka_k\Big)^{\frac{n}{n-1}}
    =CD^{\frac{n}{n-1}}\vol M\,\|dv\|_1^{\frac{n}{n-1}},
  \]
  using $\sum_ka_k=\big\||dv|\sum_k\chi_{\{2^{k-2}<v\le2^{k-1}\}}\big\|_1=\|dv\|_1$
  (the indicator functions partition the support of $v$). Hence
  \[
    \int v^{\frac{n}{n-1}}\,\vol\le CD^{\frac{n}{n-1}}\vol M\,\|dv\|_1^{\frac{n}{n-1}},
  \]
  i.e. $\|v\|_{n/(n-1)}\le CD\|dv\|_1$, as required for $v=v^+,v^-$; combining
  with $\|u-m\|_{n/(n-1)}\le\|v^+\|_{n/(n-1)}+\|v^-\|_{n/(n-1)}$ and
  $|dv^+|+|dv^-|\le|du|$ gives $\|u-u_M\|_{n/(n-1)}\le2\|u-m\|_{n/(n-1)}\le
  CD\|du\|_1$, the theorem.
\end{proof}

\begin{remark}[Remark 7.1.16]
  \label{rem:pet-ch7-dirichlet-poincare}
  \lean{PetersenLib.exerciseDirichletPoincareReference}
  \uses{thm:pet-ch7-poincare-sobolev}
  See \cref{rem:pet-ch7-ex-17} for the Poincaré inequality for functions with
  Dirichlet boundary conditions vanishing on a ball's boundary.
  \source{petersen:page-0307}
\end{remark}

\begin{proposition}[Proposition 7.1.17]
  \label{prop:pet-ch7-sobolev-hierarchy}
  \lean{PetersenLib.sobolevInterpolationHierarchy}
  \uses{thm:pet-ch7-poincare-sobolev}
  Suppose every smooth function on $(M,g)$ satisfies $\|u-u_M\|_{s-1}\le
  S\|du\|_1$ for some fixed $s>1$. Then for $1\le p<s$,
  \[
    \|u\|_{\frac{sp}{s-p}}\le\frac{p(s-1)}{s-p}S\|du\|_p+\|u\|_p.
  \]
  \source{petersen:page-0307}
\end{proposition}
\begin{proof}
  \uses{prop:pet-ch7-sobolev-hierarchy}
  For $p=1$ this follows from
  \[
    \|u-u_M\|_{s-1}\ge\|u\|_{s-1}-|u_M|\ge\|u\|_{s-1}-\|u\|_1,
  \]
  the first inequality by the reverse triangle inequality for the norm
  $\|\cdot\|_{s-1}$ and the second since $|u_M|\le\|u\|_1$.

  For $p>1$, apply the hypothesis to $u^q$ (for $q>1$ to be chosen):
  \[
    \|u\|_{\frac{q}{s-1}}^{q}=\|u^q\|_{s-1}
    \le S\|du^q\|_1+\|u^q\|_1
    =Sq\,\|u^{q-1}du\|_1+\|u^{q-1}u\|_1.
  \]
  By Hölder with exponents $p$ and $p/(p-1)$,
  \[
    \|u^{q-1}du\|_1\le\|u\|_{\frac{(q-1)p}{p-1}}^{q-1}\|du\|_p,\qquad
    \|u^{q-1}u\|_1=\|u^q\|_1\le\|u\|_{\frac{qp}{p-1}}^{q-1}\|u\|_p,
  \]
  so choosing $q=\dfrac{p(s-1)}{s-p}$, for which $\dfrac{q}{s-1}=\dfrac{sp}{s-p}$
  and $\dfrac{(q-1)p}{p-1}$ likewise resolves to $\dfrac{sp}{s-p}$, gives
  \[
    \|u\|_{\frac{sp}{s-p}}^{q}\le\|u\|_{\frac{sp}{s-p}}^{q-1}
      \Big(Sq\,\|du\|_p+\|u\|_p\Big),
  \]
  and dividing by $\|u\|_{sp/(s-p)}^{q-1}$ (when nonzero; the case
  $\|u\|_{sp/(s-p)}=0$ is trivial) gives
  \[
    \|u\|_{\frac{sp}{s-p}}\le Sq\,\|du\|_p+\|u\|_p
    =\frac{p(s-1)}{s-p}S\|du\|_p+\|u\|_p. \qedhere
  \]
\end{proof}

\begin{theorem}[Theorem 7.1.18 — Rellich Compactness]
  \label{thm:pet-ch7-rellich-compactness}
  \lean{PetersenLib.rellichCompactness}
  \uses{thm:pet-ch7-weak-poincare-sobolev-estimate}
  Assume $(M^n,g)$ is a compact Riemannian $n$-manifold. Then the inclusion
  $W^{1,2}(M)\subset L^2(M)$ (where $W^{1,2}(M)$ is the Hilbert-space closure
  of $C^\infty(M)$ in the norm $\|u\|_2^2+\|du\|_2^2$) is compact.
  \source{petersen:page-0308}
\end{theorem}
\begin{proof}
  \uses{thm:pet-ch7-rellich-compactness,thm:pet-ch7-weak-poincare-sobolev-estimate}
  Let $(u_i)$ be a sequence of smooth functions with $\|u_i\|_2^2+\|du_i\|_2^2$
  bounded. By reflexivity of the Hilbert space $W^{1,2}(M)$, a bounded
  sequence has a weakly convergent subsequence $u_i\rightharpoonup u$; in
  particular $u_{i,B(x,R)}\to u_{B(x,R)}$ for each fixed $x\in M$, $R>0$ (pairing
  weakly against the indicator of $B(x,R)$, which lies in $L^2(M)\subset
  (W^{1,2}(M))^*$ since $M$ is compact). By the Lebesgue differentiation
  theorem, $u_{B(x,R)}\to u(x)$ as $R\to0$ for almost every $x$.

  By \cref{thm:pet-ch7-weak-poincare-sobolev-estimate},
  \[
    \vol\{|u_i-u_{i,B(x,R)}|>\varepsilon\}
    \le C\Big(\frac{R}{\varepsilon}\Big)^{\frac{n}{n-1}}\vol M\,\|du_i\|_1^{\frac{n}{n-1}}
    \le C'\Big(\frac{R}{\varepsilon}\Big)^{\frac{n}{n-1}}
  \]
  with $C'$ independent of $i$ (using the uniform bound on $\|du_i\|_2$, hence
  on $\|du_i\|_1$ by Cauchy–Schwarz on the compact $M$). Letting $R\to0$ shows
  $\vol\{|u_i(x)-u(x)|>\varepsilon\}\to0$ as $i\to\infty$, for every fixed
  $\varepsilon>0$; passing to a further subsequence gives pointwise almost
  everywhere convergence $u_i\to u$. Since $\|u_i\|_{n/(n-2)}$ (for $n>2$; a
  similar bound holds for $n\le2$) is bounded by the Sobolev embedding
  following from \cref{thm:pet-ch7-poincare-sobolev} and
  \cref{prop:pet-ch7-sobolev-hierarchy}, Egorov's theorem upgrades the almost
  everywhere convergence to convergence in $L^2(M)$, as needed.
\end{proof}

\section{Applications of Ricci Curvature Comparison}

\subsection{Finiteness of Fundamental Groups}

Our first application of volume comparison controls the fundamental group. We
begin with a result on how fundamental groups can be presented.

\begin{lemma}[Lemma 7.2.1 — Gromov, 1980]
  \label{lem:pet-ch7-gromov-generators}
  \lean{PetersenLib.gromovShortGenerators}
  A compact Riemannian manifold $M$ admits generators $\{c_1,\dots,c_m\}$ for
  $\Gamma=\pi_1(M)$ such that all relations are of the form
  $c_ic_jc_k^{-1}=1$ for suitable $i,j,k$, with each generator represented by a
  loop of length $\le3\diam(M)$.
  \source{petersen:page-0308,petersen:page-0309}
\end{lemma}
\begin{proof}
  \uses{lem:pet-ch7-gromov-generators}
  For $\varepsilon\in(0,\operatorname{inj}(M))$ choose a triangulation of $M$
  whose adjacent vertices are joined by curves of length $<\varepsilon$; let
  $\{x_1,\dots,x_k\}$ be the vertices and $\{e_{ij}\}$ the edges. Fix a
  basepoint $x$, the projection of a lift $\tilde x\in\tilde M$, and join $x$
  to each $x_i$ by a segment $\sigma_i$. For adjacent vertices set
  $\sigma_{ij}=\sigma_ie_{ij}\sigma_j^{-1}$.

  Every loop based at $x$ is homotopic to a loop in the 1-skeleton of the
  triangulation, i.e.\ built from the edges $e_{ij}$; since
  $e_{ij}e_{jk}=e_{ij}\sigma_j^{-1}\sigma_je_{jk}$, such loops are products of
  loops $\sigma_{ij}$, so $\Gamma$ is generated by the $\sigma_{ij}$.

  If three vertices $x_i,x_j,x_k$ are mutually adjacent they span a $2$-simplex
  $\triangle_{ijk}$, and $\sigma_{ij}\sigma_{jk}\sigma_{ik}^{-1}$ is
  homotopically trivial (it bounds $\triangle_{ijk}$). We claim these are the
  only relations needed: if $\sigma$ is a loop in the $1$-skeleton that is
  homotopically trivial in $M$, it also contracts in the $2$-skeleton, so a
  homotopy is a union of $2$-simplices $\triangle_{ijk}$, expressing the
  relation $\sigma=1$ as a product of elementary relations
  $\sigma_{ij}\sigma_{jk}\sigma_{ik}^{-1}=1$.

  Each generator $\sigma_{ij}$ has length $\le|\sigma_i|+|e_{ij}|+|\sigma_j|\le
  2\diam(M)+\varepsilon$, which for $\varepsilon$ small is $\le3\diam(M)$,
  proving the stated bound.
\end{proof}

\begin{example}[Example 7.2.2]
  \label{ex:pet-ch7-lens-spaces-example}
  \lean{PetersenLib.exampleLensSpacesGrowingFundamentalGroup}
  \uses{lem:pet-ch7-gromov-generators}
  Consider $M_k=S^3/\mathbb{Z}_k$, the constant curvature $3$-sphere quotient by
  the cyclic group of order $k$. As $k\to\infty$, $\vol M_k\to0$ while the
  curvature stays $1$ and $\diam M_k=\pi/2$. Thus $\pi_1(M_k)=\mathbb{Z}_k$ grows
  only at the expense of shrinking volume; writing $\mathbb{Z}_k$ via generators as
  in \cref{lem:pet-ch7-gromov-generators} requires the number of generators to
  go to infinity with $k$, foreshadowing \cref{thm:pet-ch7-anderson-finiteness}.
  \source{petersen:page-0309}
\end{example}

\begin{definition}[the class $\mathfrak{M}(n,k,v,D)$]
  \label{def:pet-ch7-manifold-class}
  \lean{PetersenLib.ManifoldClass}
  For $n\in\mathbb{N}$, $k\in\mathbb{R}$, $v,D\in(0,\infty)$, let
  $\mathfrak{M}(n,k,v,D)$ denote the class of compact Riemannian $n$-manifolds
  $(M,g)$ with $\Ric\ge(n-1)k$, $\vol M\ge v$, $\diam M\le D$.
  \source{petersen:page-0309}
\end{definition}

\begin{theorem}[Theorem 7.2.3 — Anderson, 1990]
  \label{thm:pet-ch7-anderson-finiteness}
  \lean{PetersenLib.andersonFundamentalGroupFiniteness}
  \uses{def:pet-ch7-manifold-class,lem:pet-ch7-gromov-generators}
  For fixed $n,k,v,D$, there are only finitely many isomorphism types of
  fundamental group among the manifolds in $\mathfrak{M}(n,k,v,D)$.
  \source{petersen:page-0309}
\end{theorem}
\begin{proof}
  \uses{thm:pet-ch7-anderson-finiteness,lem:pet-ch7-gromov-generators,lem:pet-ch7-absolute-volume-comparison}
  Choose generators $\{c_1,\dots,c_m\}$ as in
  \cref{lem:pet-ch7-gromov-generators}. Since the number of possible relations
  among $m$ generators of the stated form is bounded by $2^m$, it suffices to
  bound $m$ in terms of $n,k,v,D$ alone. Fix $x\in M$ and view the $c_i$ as
  deck transformations of the universal cover $\tilde M\to M$; the lemma gives
  $|xc_i(x)|\le3D$. Choose a Dirichlet fundamental domain
  \[
    F=\{z\in\tilde M\mid|xz|\le|c(z)|\text{ for all }c\in\pi_1(M)\},
  \]
  a closed set with $\pi:F\to M$ onto and $\vol F=\vol M$. The translates
  $c_i(F)$ are pairwise disjoint up to measure zero, all have volume $\vol M$,
  and all lie in $B(x,6D)$ (since $|xc_i(x)|\le3D$ and $\diam F\le3D$, so
  points of $c_i(F)$ are within $3D+3D=6D$ of $x$). Hence, by absolute volume
  comparison (\cref{lem:pet-ch7-absolute-volume-comparison}),
  \[
    m\le\frac{\vol B(x,6D)}{\vol F}\le\frac{v(n,k,6D)}{v},
  \]
  a bound depending only on $n,k,v,D$.
\end{proof}

\begin{lemma}[Lemma 7.2.4 — Anderson, 1990]
  \label{lem:pet-ch7-anderson-short-loops}
  \lean{PetersenLib.andersonShortLoopSubgroupFinite}
  \uses{def:pet-ch7-manifold-class}
  For fixed $n,k,v,D$ there are $L=L(n,k,v,D)$, $N=N(n,k,v,D)$ such that for
  every $M\in\mathfrak{M}(n,k,v,D)$, any subgroup of $\pi_1(M)$ generated by
  loops of length $\le L$ has order $\le N$.
  \source{petersen:page-0310}
\end{lemma}
\begin{proof}
  \uses{lem:pet-ch7-anderson-short-loops,lem:pet-ch7-absolute-volume-comparison}
  Let $\Gamma\subset\pi_1(M)$ be generated by loops $\{c_1,\dots,c_\ell\}$ of
  length $\le L$, based at $\pi(x)$ for some lift $x\in\tilde M$, and let $F$
  be a Dirichlet fundamental domain containing $x$ as in the proof of
  \cref{thm:pet-ch7-anderson-finiteness}; distinct deck transformations move
  $F$ to sets that are disjoint up to measure zero.

  Let $U(r)\subset\Gamma$ be those elements expressible as a product of at
  most $r$ of the $c_i$'s. Since $|xc_i(x)|\le L$, every $c\in U(r)$ satisfies
  $|xc(x)|\le rL$, so $c(F)\subset B(x,rL+D)$ (as $\diam F\le D$). The sets
  $c(F)$, $c\in U(r)$, being pairwise disjoint up to measure zero and of
  volume $\vol F\ge v$, absolute volume comparison
  (\cref{lem:pet-ch7-absolute-volume-comparison}) gives
  \[
    |U(r)|\le\frac{\vol B(x,rL+D)}{\vol F}\le\frac{v(n,k,rL+D)}{v}.
  \]
  Set $N=v(n,k,2D)/v+1$ and $L=D/N$. If $\Gamma$ had more than $N$ elements,
  taking $r=N$ would give $rL+D=D+D=2D$, so
  \[
    N+1\le|U(N)|+1\le\Gamma
    \quad\text{forces}\quad
    N< |U(N)|\le\frac{v(n,k,2D)}{v}=N-1,
  \]
  a contradiction (more precisely: if $|\Gamma|>N$ then some element of
  $\Gamma$ lies outside $U(N)$ only if $U(N)\ne\Gamma$, but
  $|U(N)|\le v(n,k,2D)/v=N-1<N<|\Gamma|$ already contradicts $U(N)=\Gamma$
  when $\Gamma$ is generated by the $c_i$ and has more than $N$ elements
  reachable within $N$ products). Hence $|\Gamma|\le N$.
\end{proof}

\subsection{Maximal Diameter Rigidity}

Given Myers' diameter estimate, it is natural to ask what happens when the
diameter attains its maximal value; only the sphere has this property.

\begin{theorem}[Theorem 7.2.5 — S. Y. Cheng, 1975]
  \label{thm:pet-ch7-cheng-rigidity}
  \lean{PetersenLib.chengMaximalDiameterRigidity}
  \uses{lem:pet-ch7-laplacian-comparison-calabi,thm:pet-ch7-strong-maximum-principle}
  If $(M,g)$ is complete with $\Ric\ge(n-1)k>0$ and $\diam=\pi/\sqrt k$, then
  $(M,g)$ is isometric to $S^n_k$ (the space form of curvature $k$).
  \source{petersen:page-0310,petersen:page-0311,petersen:page-0312}
\end{theorem}
\begin{proof}
  \uses{thm:pet-ch7-cheng-rigidity,lem:pet-ch7-laplacian-comparison-calabi,thm:pet-ch7-strong-maximum-principle,prop:pet-ch7-volume-form-laplacian-trace}
  Fix $p,q\in M$ with $|pq|=\pi/\sqrt k$ and set $r(x)=|xp|$,
  $\tilde r(x)=|xq|$. We establish successively:
  \begin{enumerate}
    \item[(1)] $r+\tilde r=\pi/\sqrt k$ on $M$.
    \item[(2)] $r,\tilde r$ are smooth on $M-\{p,q\}$.
    \item[(3)] $\Hess r=\dfrac{\snk_k'(r)}{\snk_k(r)}\,g_r$ on $M-\{p,q\}$,
      where $g_r$ is the metric restricted to the distribution orthogonal to
      $\partial_r$.
    \item[(4)] $g=dr^2+\snk_k^2(r)\,ds^2_{n-1}$, which forces $M\cong S^n_k$.
  \end{enumerate}

  \emph{Proof of (1).} The triangle inequality gives $r+\tilde r\ge\pi/\sqrt k$
  everywhere, with equality on any segment joining $p,q$. By
  \cref{lem:pet-ch7-laplacian-comparison-calabi} applied at $p$ and at $q$
  (with $\snk_k'(t)/\snk_k(t)=\sqrt k\cot(\sqrt kt)$ for $k>0$),
  \[
    \Delta(r+\tilde r)\le(n-1)\sqrt k\cot(\sqrt kr)+(n-1)\sqrt k\cot(\sqrt k\tilde r)
    \le(n-1)\sqrt k\Big(\cot(\sqrt kr)+\cot\big(\pi-\sqrt kr\big)\Big)=0,
  \]
  using $\tilde r\ge\pi/\sqrt k-r$ and monotonicity of $\cot$ on $(0,\pi)$ to
  replace $\tilde r$ by $\pi/\sqrt k-r$ in the second term, then
  $\cot(\pi-\theta)=-\cot\theta$. Thus $r+\tilde r$ is superharmonic on
  $M-\{p,q\}$ and attains its global minimum $\pi/\sqrt k$ on the segment; the
  minimum principle (\cref{thm:pet-ch7-strong-maximum-principle} applied to
  $-(r+\tilde r)$) forces $r+\tilde r\equiv\pi/\sqrt k$ on $M$.

  \emph{Proof of (2).} If $x\in M-\{p,q\}$, join $x$ to $p$ and to $q$ by
  segments $c_1,c_2$. By (1), concatenating $c_1$ and $c_2$ gives a curve from
  $p$ to $q$ of length $\pi/\sqrt k=|pq|$, hence a segment, which must be
  smooth; so $c_1,c_2$ are subsegments of a single segment through $x$. The
  standard characterization of smoothness of distance functions (via absence
  of conjugate/cut points) then gives smoothness of $r,\tilde r$ at $x$.

  \emph{Proof of (3).} Since $r+\tilde r\equiv\pi/\sqrt k$, $\Delta
  r=-\Delta\tilde r$. Applying \cref{lem:pet-ch7-laplacian-comparison-calabi}
  to both $r$ and $\tilde r=\pi/\sqrt k-r$,
  \[
    (n-1)\frac{\snk_k'(r)}{\snk_k(r)}
    \ge\Delta r=-\Delta\tilde r
    \ge-(n-1)\frac{\snk_k'(\tilde r)}{\snk_k(\tilde r)}
    =-(n-1)\frac{\snk_k'(\pi/\sqrt k-r)}{\snk_k(\pi/\sqrt k-r)}
    =(n-1)\frac{\snk_k'(r)}{\snk_k(r)},
  \]
  the last equality since $\snk_k(\pi/\sqrt k-r)=\snk_k(r)$ and
  $\snk_k'(\pi/\sqrt k-r)=-\snk_k'(r)$ (symmetry of $\sin(\sqrt k\,\cdot)$
  about $\pi/(2\sqrt k)$ for $k>0$). Hence equality holds throughout:
  $\Delta r=(n-1)\snk_k'(r)/\snk_k(r)$, and moreover, since the chain of
  inequalities of (tr2) from \cref{prop:pet-ch7-volume-form-laplacian-trace}
  (equivalently, the Bochner/Riccati inequality underlying
  \cref{lem:pet-ch7-laplacian-comparison-calabi})
  \[
    -(n-1)k=\partial_r(\Delta r)+\frac{(\Delta r)^2}{n-1}\le\partial_r(\Delta
    r)+|\Hess r|^2\le-\Ric(\partial_r,\partial_r)\le-(n-1)k
  \]
  is now an equality throughout, we get in particular $(\Delta
  r)^2=(n-1)|\Hess r|^2$ — precisely the equality case, recalled from the
  proof of \cref{prop:pet-ch7-volume-form-laplacian-trace}, of the
  Cauchy--Schwarz inequality $(\operatorname{tr}A)^2\le(n-1)|A|^2$ applied to
  $A=\Hess r$ restricted to
  the $(n-1)$-dimensional distribution orthogonal to $\partial_r$ (on which
  $\Hess r(\partial_r,\cdot)=0$ since $|\nabla r|\equiv1$). Equality forces
  $A$ to be a scalar multiple of the identity on that distribution:
  \[
    \Hess r=\frac{\Delta r}{n-1}\,g_r=\frac{\snk_k'(r)}{\snk_k(r)}\,g_r.
  \]

  \emph{Proof that (3) implies (4), and that (4) forces $M\cong S^n_k$.} These
  are established by the general theory of rotationally symmetric metrics
  determined by their radial Hessian (the ODE $\Hess
  r=(\snk_k'/\snk_k)g_r$ integrates, in polar coordinates about $p$, exactly
  to the metric $dr^2+\snk_k^2(r)ds^2_{n-1}$ on $M-\{p,q\}$, which extends
  smoothly across $p$ and $q$ by the smoothness criterion for rotationally
  symmetric metrics), giving the isometry with the space form $S^n_k$.
\end{proof}

\begin{example}[Example 7.2.6 — Anderson, 1990]
  \label{ex:pet-ch7-anderson-cp2-example}
  \lean{PetersenLib.andersonAlmostMaximalDiameterExampleFour}
  \uses{thm:pet-ch7-cheng-rigidity}
  For $n=4$, consider metrics $dr^2+\rho^2\sigma_1^2+\phi^2(\sigma_2^2+\sigma_3^2)$
  on $I\times S^3$, with
  \[
    \rho(r)=\begin{cases}\dfrac{\sin(ar)}{a}&r\le r_0,\\c_1\sin(r+\delta)&r\ge
    r_0,\end{cases}\qquad
    \phi(r)=\begin{cases}br^2+c&r\le r_0,\\c_2\sin(r+\delta)&r\ge
    r_0,\end{cases}
  \]
  reflected in $r=\pi/2-\delta$ to produce a metric on $\mathbb{CP}^2\#\overline{\mathbb{CP}^2}$.
  For any small $r_0>0$ the parameters can be tuned so $\rho,\phi$ are $C^1$
  and generate a metric with $\Ric\ge3$; as $r_0\to0$ the parameter interval
  $I\to[0,\pi]$, so the diameters converge to $\pi$, showing that
  \cref{thm:pet-ch7-cheng-rigidity} does not admit a topological rigidity
  refinement when $n=4$.
  \source{petersen:page-0312,petersen:page-0313}
\end{example}

\begin{example}[Example 7.2.7 — Otsu, 1991]
  \label{ex:pet-ch7-otsu-example}
  \lean{PetersenLib.otsuAlmostMaximalDiameterExample}
  \uses{ex:pet-ch7-anderson-cp2-example}
  For $n\ge5$, standard doubly warped products $dr^2+\rho^2ds_2^2+\phi^2ds^2_{n-3}$
  on $I\times S^2\times S^{n-3}$, with $\rho,\phi$ chosen analogously to
  \cref{ex:pet-ch7-anderson-cp2-example}, yield metrics on $S^2\times
  S^{n-2}$ with $\Ric\ge n-1$ and diameter $\to\pi$. In both examples the
  functions $\rho,\phi$ are only $C^1$ at the outset, but being concave they
  can be smoothed near the break points while remaining concave, changing
  values and first derivatives negligibly and only increasing the (already
  nonpositive) second derivative in absolute value, so the lower Ricci bound
  persists after smoothing.
  \source{petersen:page-0313}
\end{example}

\section{Manifolds of Nonnegative Ricci Curvature}

We prove the splitting theorem of Cheeger and Gromoll, analogous in many ways
to the maximal diameter theorem, with far-reaching consequences for compact
manifolds of nonnegative Ricci curvature — for instance it shows that
$S^3\times S^1$ admits no Ricci flat metric. Throughout this section $(M,g)$
is complete and noncompact.

\subsection{Rays and Lines}

\begin{definition}[ray, line]
  \label{def:pet-ch7-ray-line}
  \lean{PetersenLib.IsRay,PetersenLib.IsLine}
  A \emph{ray} $r:[0,\infty)\to(M,g)$ is a unit-speed geodesic with
  $|r(t)r(s)|=|t-s|$ for all $t,s\ge0$; it may be viewed as a semi-infinite
  segment, or a segment from $r(0)$ to infinity. A \emph{line}
  $l:\mathbb{R}\to(M,g)$ is a unit-speed geodesic with $|l(t)l(s)|=|t-s|$ for all
  $t,s\in\mathbb{R}$.
  \source{petersen:page-0313}
\end{definition}

\begin{definition}[connected/disconnected at infinity]
  \label{def:pet-ch7-connected-infinity}
  \lean{PetersenLib.IsConnectedAtInfinity}
  A complete manifold $(M,g)$ is \emph{connected at infinity} if for every
  compact $K\subset M$ there is a compact $C\supset K$ such that any two
  points of $M-C$ can be joined by a curve in $M-K$. Otherwise $(M,g)$ is
  \emph{disconnected at infinity}.
  \source{petersen:page-0314}
\end{definition}

\begin{lemma}[Lemma 7.3.1]
  \label{lem:pet-ch7-ray-line-existence}
  \lean{PetersenLib.rayExistenceLineIfDisconnectedAtInfinity}
  \uses{def:pet-ch7-ray-line,def:pet-ch7-connected-infinity}
  If $p\in(M,g)$ then a ray emanating from $p$ always exists. If $M$ is
  disconnected at infinity, $(M,g)$ contains a line.
  \source{petersen:page-0313,petersen:page-0314}
\end{lemma}
\begin{proof}
  \uses{lem:pet-ch7-ray-line-existence}
  Fix $p\in M$ and a sequence $q_i\to\infty$; let $v_i\in T_pM$ be unit
  vectors with $\sigma_i(t)=\exp_p(tv_i)$, $t\in[0,d(p,q_i)]$, a segment from
  $p$ to $q_i$. Passing to a subsequence, $v_i\to v\in T_pM$. Then
  $\sigma(t)=\exp_p(tv)$, $t\in[0,\infty)$, is a segment: by continuity of
  $\exp_p$, $\sigma_i\to\sigma$ pointwise, so
  $|\sigma(s)\sigma(t)|=\lim_i|\sigma_i(s)\sigma_i(t)|=|s-t|$ for all
  $s,t\ge0$ where both sides are eventually defined (using $d(p,q_i)\to\infty$).
  This $\sigma$ is the desired ray.

  If $M$ is disconnected at infinity, there is a compact $K$ and sequences
  $p_i,q_i\to\infty$ such that every curve from $p_i$ to $q_i$ meets $K$.
  Joining $p_i,q_i$ by segments $\sigma_i:(-a_i,b_i)\to M$ with $a_i,b_i\to
  \infty$ and $\sigma_i(0)\in K$ (possible since $\sigma_i$ must cross $K$, and
  reparametrizing so the crossing time is $0$), the same subconvergence
  argument as above (using compactness of $K$ to extract a convergent
  subsequence of basepoints and directions) produces a line
  $\sigma(t)=\lim_i\sigma_i(t)$, $t\in\mathbb{R}$.
\end{proof}

\begin{example}[Example 7.3.2]
  \label{ex:pet-ch7-paraboloid-no-lines}
  \lean{PetersenLib.paraboloidContainsNoLines}
  \uses{def:pet-ch7-ray-line}
  Surfaces of revolution $dr^2+\rho^2(r)ds^2_{n-1}$ with $\rho:[0,\infty)\to
  [0,\infty)$, $\dot\rho(t)<1$, $\ddot\rho(t)<0$ for $t>0$, contain no lines;
  such manifolds look like paraboloids.
  \source{petersen:page-0314}
\end{example}

\begin{example}[Example 7.3.3]
  \label{ex:pet-ch7-cylinder-contains-line}
  \lean{PetersenLib.sphereCrossRAlwaysContainsLine}
  \uses{lem:pet-ch7-ray-line-existence}
  Any complete metric on $S^{n-1}\times\mathbb{R}$ contains a line, since the
  manifold is disconnected at infinity.
  \source{petersen:page-0315}
\end{example}

\begin{example}[Example 7.3.4]
  \label{ex:pet-ch7-schwarzschild-no-line}
  \lean{PetersenLib.schwarzschildContainsNoLine}
  \uses{def:pet-ch7-ray-line}
  The Schwarzschild metric on $S^{n-2}\times\mathbb{R}^2$ contains no lines, as
  follows from the splitting theorem (\cref{thm:pet-ch7-splitting-theorem})
  since the space is not metrically a product.
  \source{petersen:page-0315}
\end{example}

\subsection{Busemann Functions}

For the remainder of §7.3, $(M,g)$ is a complete noncompact manifold with
$\Ric\ge0$. Given a unit-speed ray $c:[0,\infty)\to(M,g)$, define
$b_t(x)=|xc(t)|-t$.

\begin{proposition}[Proposition 7.3.6]
  \label{prop:pet-ch7-busemann-approximants}
  \lean{PetersenLib.busemannApproximantsProperties}
  \uses{def:pet-ch7-ray-line,lem:pet-ch7-laplacian-comparison-calabi}
  The functions $b_t$, $t\ge0$, satisfy:
  \begin{enumerate}
    \item[(1)] for fixed $x$, $t\mapsto b_t(x)$ is decreasing and bounded in
      absolute value by $|xc(0)|$;
    \item[(2)] $|b_t(x)-b_t(y)|\le|xy|$;
    \item[(3)] $\Delta b_t(x)\le\dfrac{n-1}{b_t(x)+t}$ everywhere.
  \end{enumerate}
  \source{petersen:page-0316}
\end{proposition}
\begin{proof}
  \uses{prop:pet-ch7-busemann-approximants}
  (2) is the triangle inequality applied to $|xc(t)|,|yc(t)|,|xy|$, giving
  $|b_t(x)-b_t(y)|=\big||xc(t)|-|yc(t)|\big|\le|xy|$. (3) follows since
  $b_t(x)+t=|xc(t)|$ is (up to the additive constant $t$) the distance
  function from $c(t)$, to which \cref{lem:pet-ch7-laplacian-comparison-calabi}
  applies with $k=0$: $\Delta(b_t+t)(x)\le(n-1)\snk_0'(r)/\snk_0(r)=(n-1)/r$
  with $r=|xc(t)|=b_t(x)+t$.

  For (1), the triangle inequality gives
  \[
    |b_t(x)|=\big||xc(t)|-|c(0)c(t)|\big|\le|xc(0)|.
  \]
  For $s<t$,
  \[
    b_t(x)-b_s(x)=|xc(t)|-t-|xc(s)|+s
    =|xc(t)|-|xc(s)|-|c(t)c(s)|
    \le|c(t)c(s)|-|c(t)c(s)|=0,
  \]
  using $|xc(t)|-|xc(s)|\le|c(t)c(s)|$ (triangle inequality) and
  $t-s=|c(t)c(s)|$ (since $c$ is a ray). Thus $t\mapsto b_t(x)$ is
  nonincreasing, hence decreasing (as an infimum bounded below by
  $-|xc(0)|$).
\end{proof}

\begin{definition}[Busemann function]
  \label{def:pet-ch7-busemann-function}
  \lean{PetersenLib.busemannFunction}
  \uses{prop:pet-ch7-busemann-approximants}
  By \cref{prop:pet-ch7-busemann-approximants}, $b_t$ converges pointwise as
  $t\to\infty$ to a $1$-Lipschitz function
  \[
    b_\infty(x)=\lim_{t\to\infty}b_t(x),\qquad |b_\infty(x)-b_\infty(y)|\le|xy|,
    \qquad |b_\infty(x)|\le|xc(0)|,
  \]
  satisfying $b_\infty(c(r))=\lim_{t\to\infty}(|c(r)c(t)|-t)=-r$. This
  $b_\infty$, denoted $b_c$, is the \emph{Busemann function} for $c$, thought
  of as a renormalized distance function from ``$c(\infty)$.''
  \source{petersen:page-0316}
\end{definition}

\begin{example}[Example 7.3.7]
  \label{ex:pet-ch7-euclidean-busemann}
  \lean{PetersenLib.euclideanBusemannFunction}
  \uses{def:pet-ch7-busemann-function}
  If $M=(\mathbb{R}^n,g_{\mathbb{R}^n})$, every Busemann function has the form
  $b_c(x)=\dot c(0)\cdot(c(0)-x)$.
  \source{petersen:page-0316}
\end{example}

\begin{definition}[horosphere, asymptote]
  \label{def:pet-ch7-horosphere-asymptote}
  \lean{PetersenLib.Horosphere,PetersenLib.IsAsymptoteFor}
  \uses{def:pet-ch7-busemann-function}
  The level sets $b_c^{-1}(t)$ are \emph{horospheres} (hyperplanes in $\mathbb{R}^n$;
  spheres tangent to the boundary in the Poincaré model of hyperbolic space).
  Given a ray $c$ and $p\in M$, a family of unit-speed segments
  $\sigma_i:[0,L_i]\to M$ from $p$ to $c(t_i)$, $t_i\to\infty$, subconverges (as
  in \cref{lem:pet-ch7-ray-line-existence}) to a ray $\tilde c:[0,\infty)\to
  M$, $\tilde c(0)=p$, called an \emph{asymptote} for $c$ from $p$; asymptotes
  need not be unique.
  \source{petersen:page-0317}
\end{definition}

\begin{proposition}[Proposition 7.3.8]
  \label{prop:pet-ch7-busemann-asymptote-relation}
  \lean{PetersenLib.busemannAsymptoteRelation}
  \uses{def:pet-ch7-horosphere-asymptote}
  With $c,p,\tilde c$ as in \cref{def:pet-ch7-horosphere-asymptote}, the
  Busemann functions satisfy
  \begin{enumerate}
    \item[(1)] $b_c(x)\le b_c(p)+b_{\tilde c}(x)$ for all $x\in M$;
    \item[(2)] $b_c(\tilde c(t))=b_c(p)+b_{\tilde c}(\tilde c(t))=b_c(p)-t$ for
      all $t\ge0$.
  \end{enumerate}
  \source{petersen:page-0317,petersen:page-0318}
\end{proposition}
\begin{proof}
  \uses{prop:pet-ch7-busemann-asymptote-relation}
  Let $\sigma_i:[0,L_i]\to M$ be segments from $p$ to $c(t_i)$ converging to
  $\tilde c$. For (2): since $\sigma_i$ is a segment,
  $|pc(t_i)|=|p\sigma_i(s)|+|\sigma_i(s)c(t_i)|$ for any fixed $s\le L_i$, and
  $\sigma_i(s)\to\tilde c(s)$ as $i\to\infty$. Hence
  \[
    b_c(p)=\lim_i(|pc(t_i)|-t_i)
    =\lim_i(|p\sigma_i(s)|+|\sigma_i(s)c(t_i)|-t_i)
    =|p\tilde c(s)|+\lim_i(|\tilde c(s)c(t_i)|-t_i)
    =s+b_c(\tilde c(s)),
  \]
  i.e.\ $b_c(p)-s=b_c(\tilde c(s))$; since $\tilde c$ is unit speed with
  $\tilde c(0)=p$, $b_{\tilde c}(\tilde c(s))=-s$
  (\cref{def:pet-ch7-busemann-function}), so $b_c(p)+b_{\tilde
  c}(\tilde c(s))=b_c(p)-s=b_c(\tilde c(s))$, proving (2).

  For (1), fix $x\in M$ and $t\ge0$; using the triangle inequality through
  $\tilde c(t)$ and then through $p$,
  \[
    |xc(s)|-s\le|x\tilde c(t)|+|\tilde c(t)c(s)|-s.
  \]
  Letting $s\to\infty$ and using $b_c(x)=\lim_s(|xc(s)|-s)$ on the left and
  $\lim_s(|\tilde c(t)c(s)|-s)=b_c(\tilde c(t))$ on the right,
  \[
    b_c(x)\le|x\tilde c(t)|+b_c(\tilde c(t))
    =|x\tilde c(t)|-t+t+b_c(\tilde c(t)).
  \]
  As $t\to\infty$, $|x\tilde c(t)|-t\to b_{\tilde c}(x)$, and by (2)
  $b_c(\tilde c(t))+t=b_c(p)$ is constant in $t$, giving
  $b_c(x)\le b_{\tilde c}(x)+b_c(p)$, which is (1).
\end{proof}

\begin{lemma}[Lemma 7.3.9]
  \label{lem:pet-ch7-busemann-superharmonic}
  \lean{PetersenLib.busemannFunctionSuperharmonic}
  \uses{prop:pet-ch7-busemann-asymptote-relation,prop:pet-ch7-busemann-approximants}
  If $\Ric(M,g)\ge0$, then $\Delta b_c\le0$ everywhere.
  \source{petersen:page-0318}
\end{lemma}
\begin{proof}
  \uses{lem:pet-ch7-busemann-superharmonic}
  Fix $p\in M$ and an asymptote $\tilde c$ for $c$ from $p$. By
  \cref{prop:pet-ch7-busemann-asymptote-relation}(1), $b_c(p)+b_{\tilde c}$ is
  a support function from above for $b_c$ at $p$ (they agree at $p$ by (2),
  and dominate $b_c$ elsewhere by (1)), so it suffices to show $\Delta
  b_{\tilde c}\le0$ at $p$. The functions $b_t(x)=|x\tilde c(t)|-t$ are
  themselves support functions from above for $b_{\tilde c}$ at $p$ (by
  \cref{prop:pet-ch7-busemann-approximants}(1), since $b_{\tilde c}=\lim_tb_t$
  is a decreasing limit and $b_{\tilde c}(p)=\lim_tb_t(p)$, while for general
  $x$ near $p$, $b_t(x)\ge b_{\tilde c}(x)$ since the sequence $b_t(x)$
  decreases to $b_{\tilde c}(x)$). Each $b_t$ is smooth at $p$ (being, up to a
  constant, the distance function from $\tilde c(t)\ne p$) with
  \[
    \Delta b_t(p)\le\frac{n-1}{t}\to0
  \]
  as $t\to\infty$, by \cref{prop:pet-ch7-busemann-approximants}(3). Hence
  $\Delta b_{\tilde c}(p)\le0$ in the support sense, giving $\Delta
  b_c(p)\le0$.
\end{proof}

\begin{theorem}[Theorem 7.3.5 — The Splitting Theorem, Cheeger and Gromoll, 1971]
  \label{thm:pet-ch7-splitting-theorem}
  \lean{PetersenLib.cheegerGromollSplittingTheorem}
  \uses{lem:pet-ch7-busemann-superharmonic,thm:pet-ch7-strong-maximum-principle,thm:pet-ch7-harmonic-regularity}
  If $(M,g)$ contains a line and has $\Ric\ge0$, then $(M,g)$ is isometric to
  a product $(H\times\mathbb{R},\,g_0+dt^2)$.
  \source{petersen:page-0315,petersen:page-0318,petersen:page-0319}
\end{theorem}
\begin{proof}
  \uses{thm:pet-ch7-splitting-theorem,lem:pet-ch7-busemann-superharmonic,thm:pet-ch7-strong-maximum-principle,thm:pet-ch7-harmonic-regularity,prop:pet-ch7-busemann-approximants,def:pet-ch7-busemann-function,prop:pet-ch7-volume-form-laplacian-trace}
  The strategy is to find a distance function $r:M\to\mathbb{R}$ (i.e.\
  $|\nabla r|\equiv1$) that is linear ($\Hess r=0$); then $M=U_0\times\mathbb{R}$
  with $U_0=\{r=0\}$ and $g=g_0+dt^2$, as in the maximal diameter theorem
  (\cref{thm:pet-ch7-cheng-rigidity}), where two distance functions at maximal
  distance were used and shown to sum to a constant, hence (via the maximum
  principle) each be smooth away from two points with an explicit Hessian
  determined by the rigidity case of Cauchy--Schwarz. Here the two ends of a
  line play the role of the two maximally-distant points.

  Let $c:\mathbb{R}\to M$ be a line, and let $b^+,b^-$ be the Busemann functions
  for the rays $c|_{[0,\infty)}$ and $t\mapsto c(-t)$ respectively:
  \[
    b^+(x)=\lim_{t\to\infty}(|xc(t)|-t),\qquad
    b^-(x)=\lim_{t\to\infty}(|xc(-t)|-t).
  \]
  By the triangle inequality, $b^+(x)+b^-(x)=\lim_{t\to\infty}(|xc(t)|+|xc(-t)|-2t)\ge0$
  for all $x$ (each term $|xc(t)|+|xc(-t)|-2t\ge|c(t)c(-t)|-2t=0$ since $c$ is
  a line), with equality on $c(\mathbb{R})$ since $c$ realizes the distance
  $|c(t)c(-t)|=2t$. By \cref{lem:pet-ch7-busemann-superharmonic},
  $\Delta(b^++b^-)\le0$, so $b^++b^-$ is superharmonic with a global minimum
  (value $0$, attained on the line $c$). By
  \cref{thm:pet-ch7-strong-maximum-principle} applied to $-(b^++b^-)$,
  $b^++b^-\equiv0$ on $M$. Hence $b^+=-b^-$ and, since
  $\Delta b^+\le0\le-\Delta b^-=\Delta b^+$ (using $\Delta b^-\le0$ from
  \cref{lem:pet-ch7-busemann-superharmonic} applied to the opposite ray), we
  get $\Delta b^+=\Delta b^-=0$ everywhere.

  Next, $b^\pm$ have unit gradient. Let $p\in M$ and let $\tilde c^\pm$ be
  asymptotes for $c^\pm=c|_{\pm[0,\infty)}$ from $p$. Set
  $b_t^{\pm,\pm}(x)=|x\tilde c^\pm(t)|-t$; by the same support-function
  argument used to prove \cref{lem:pet-ch7-busemann-superharmonic},
  \[
    b^+(x)\ge b^+(x)-b^+(p)=-b^-(x)+b^-(p)\ge-b_t^{-,-}(x)+b_t^{-,-}(p)
  \]
  wherever these support relations hold, with equality at $x=p$; since
  $b_t^{\pm,\pm}$ are smooth at $p$ with unit gradient (being, up to sign and
  additive constant, distance functions from points $\ne p$) and $b^+,b^-$
  agree with them to first order at $p$ from opposite sides, $\nabla
  b^+(p)=-\nabla b^-(p)$ and both $b^\pm$ are differentiable at $p$ with unit
  gradient. As $p$ was arbitrary, $b^\pm$ are everywhere differentiable with
  $|\nabla b^\pm|\equiv1$, and being harmonic in the weak sense, are smooth by
  \cref{thm:pet-ch7-harmonic-regularity}.

  Finally, apply (tr2) of \cref{prop:pet-ch7-volume-form-laplacian-trace} to
  $r=b^+$ (the same rigidity argument used in
  \cref{thm:pet-ch7-cheng-rigidity}(3)), with $k=0$ and $\Delta
  r\equiv0$:
  \[
    0=\partial_r\Delta b^++\frac{(\Delta b^+)^2}{n-1}
    \le\partial_r\Delta b^++|\Hess b^+|^2
    =|\Hess b^+|^2\le-\Ric(\nabla b^+,\nabla b^+)\le0,
  \]
  forcing $|\Hess b^+|^2=0$, i.e.\ $\Hess b^+\equiv0$: $b^+$ is the sought
  linear distance function, and $M=\{b^+=0\}\times\mathbb{R}$ isometrically, with
  the $\mathbb{R}$ factor parametrized by $b^+$.
\end{proof}

\subsection{Structure Results in Nonnegative Ricci Curvature}

The splitting theorem gives several structural consequences for compact
manifolds of nonnegative Ricci curvature.

\begin{corollary}[Corollary 7.3.10]
  \label{cor:pet-ch7-no-ricci-flat-sphere-product}
  \lean{PetersenLib.noRicciFlatMetricOnSphereProductLowDim}
  \uses{thm:pet-ch7-splitting-theorem}
  $S^k\times S^l$ admits no Ricci flat metric when $k=2,3$.
  \source{petersen:page-0319}
\end{corollary}
\begin{proof}
  \uses{cor:pet-ch7-no-ricci-flat-sphere-product}
  The universal cover of $S^k\times S^l$ is $S^k\times\mathbb{R}^l\times(\text{a
  further cover of }S^l)$; take instead the cover $S^k\times\mathbb{R}$ obtained by
  unrolling one $S^1\subset S^l$ factor when $l\ge1$ (or directly $S^k\times
  \mathbb{R}$ when $l=1$). This space is disconnected at infinity, so any metric of
  nonnegative Ricci curvature on it must split
  (\cref{thm:pet-ch7-splitting-theorem}, applied after
  \cref{lem:pet-ch7-ray-line-existence} produces a line). If the original
  metric on $S^k\times S^l$ is Ricci flat, the splitting produces a Ricci
  flat metric on a compact, simply connected $k$-manifold $H$ homotopy
  equivalent to $S^k$. For $k\le3$ a Ricci flat metric is flat, but a compact
  flat manifold has infinite (virtually abelian) fundamental group by
  Bieberbach theory, contradicting simple connectivity unless $H$ is a point —
  impossible since $H$ is homotopy equivalent to $S^k$, $k\ge2$.
\end{proof}

\begin{theorem}[Theorem 7.3.11 — Structure Theorem for Nonnegative Ricci Curvature, Cheeger and Gromoll, 1971]
  \label{thm:pet-ch7-structure-theorem}
  \lean{PetersenLib.nonnegativeRicciStructureTheorem}
  \uses{thm:pet-ch7-splitting-theorem,lem:pet-ch7-ray-line-existence}
  Suppose $(M,g)$ is compact with $\Ric\ge0$. Then:
  \begin{enumerate}
    \item[(1)] the universal cover $(\tilde M,\tilde g)$ splits isometrically
      as $N\times\mathbb{R}^k$, with $N$ compact and containing no lines;
    \item[(2)] $\Iso(\tilde M)=\Iso(N)\times\Iso(\mathbb{R}^k)$;
    \item[(3)] there is a finite normal subgroup $G\subset\pi_1(M)$ with
      $\pi_1(M)/G\cong\pi_1(M)\cap\Iso(\mathbb{R}^k)$, and a finite-index subgroup
      $\mathbb{Z}^k\subset\pi_1(M)\cap\Iso(\mathbb{R}^k)$.
  \end{enumerate}
  \source{petersen:page-0319,petersen:page-0320}
\end{theorem}
\begin{proof}
  \uses{thm:pet-ch7-structure-theorem}
  \emph{(1).} Applying \cref{thm:pet-ch7-splitting-theorem} repeatedly as long
  as lines exist writes $\tilde M=N\times\mathbb{R}^k$ with $N$ containing no lines
  (the process terminates since $\dim\le n$). If $c(t)=(c_1(t),c_2(t))\in
  N\times\mathbb{R}^k$ is a geodesic, both $c_i$ are geodesics, and if $c$ is a line
  then each nonconstant $c_i$ is a line; since $N$ has no lines, $c_1$ is
  constant, so every line in $\tilde M$ has the form $c(t)=(x,\sigma(t))$
  with $x\in N$, $\sigma$ a line in $\mathbb{R}^k$.

  \emph{(2).} Let $F\in\Iso(\tilde M)$. Since every tangent vector to $\mathbb{R}^k$
  is the velocity of a line in $\{x\}\times\mathbb{R}^k$ and $F$ carries lines to
  lines, and (1) shows every line lies in some $\{x\}\times\mathbb{R}^k$, $F$
  carries $\{x\}\times\mathbb{R}^k$ to $\{F_1(x)\}\times\mathbb{R}^k$ for a well-defined
  map $F_1:N\to N$. Thus $F=(F_1,F_2)$ with $F_1:N\to N$ an isometry; since
  $DF$ preserves the tangent spaces to the $\mathbb{R}^k$ factors it also preserves
  the orthogonal complements (the tangent spaces to $N$), so $F_2:\mathbb{R}^k\to
  \mathbb{R}^k$ is an isometry. Hence $\Iso(\tilde M)=\Iso(N)\times\Iso(\mathbb{R}^k)$.

  \emph{(3), part on $N$ compact (used implicitly in (1)):} $\pi_1(M)$ acts on
  $\tilde M$ discretely and cocompactly by deck transformations; projecting to
  $\Iso(N)$ via (2), the subgroup $\pi_1(M)\cap\Iso(N)$ acts cocompactly on
  $N$ (for any sequence $p_i\in N$ one can find deck transformations $F_i$
  with $F_i(p_i)$ in a fixed compact set, since the action on $\tilde M$ is
  cocompact). If $N$ were noncompact it would contain a ray $c$; choosing
  $t_i\to\infty$ and $F_i\in\pi_1(M)\cap\Iso(N)$ with $F_i(c(t_i))$ in a
  compact set, a subsequential limit of $DF_i(\dot c(t_i))$ gives a unit
  vector $v$ such that the geodesics $c_i(t)=F_i(c(t+t_i))$ converge to
  $\exp(tv)$; since each $c_i$ is a ray on $[a,\infty)$ once $t_i\ge-a$, the
  limit $\exp(tv)$ is a ray on every $[a,\infty)$, i.e.\ a line in $N$ — a
  contradiction. Hence $N$ is compact.

  \emph{(3).} Let $G=\ker\big(\pi_1(M)\to\pi_1(M)\cap\Iso(\mathbb{R}^k)\big)$, the
  kernel of the projection to the $\mathbb{R}^k$-factor of the isometry (using (2)).
  Then $G$ acts freely and discretely on $\tilde M$ without moving the
  $\mathbb{R}^k$-factor, hence freely and discretely on the compact $N$, so $G$ is
  finite; by construction $\pi_1(M)/G\cong\pi_1(M)\cap\Iso(\mathbb{R}^k)$.

  The translation subgroup $\mathbb{R}^k\subset\Iso(\mathbb{R}^k)$ is normal with factor
  group $O(k)$. The intersection $\Lambda=\pi_1(M)\cap\mathbb{R}^k$ is a finitely
  generated abelian group acting discretely and cocompactly on $\mathbb{R}^k$ (it
  has finite index in $\pi_1(M)\cap\Iso(\mathbb{R}^k)$, which itself acts
  cocompactly since $\pi_1(M)$ acts cocompactly on $\tilde M=N\times\mathbb{R}^k$
  and $N$ is compact), so $\Lambda\cong\mathbb{Z}^m$ for some $m$. If $m<k$, $\Lambda$
  spans a proper subspace of $\mathbb{R}^k$ and cannot act cocompactly. If $m>k$,
  $\Lambda$ contains two elements linearly independent over $\mathbb{Q}$ but not
  over $\mathbb{R}$; the subgroup they generate has orbits dense in a line, so
  cannot act discretely — a contradiction. Hence $m=k$ and $\Lambda\cong\mathbb{Z}^k$
  has finite index in $\pi_1(M)\cap\Iso(\mathbb{R}^k)$.
\end{proof}

\begin{remark}[Remark 7.3.12]
  \label{rem:pet-ch7-wilking-remark}
  \lean{PetersenLib.wilkingConverseRemark}
  \uses{thm:pet-ch7-structure-theorem}
  Wilking has shown conversely that any group $G$ admitting a finite normal
  subgroup $H\subset G$ with $G/H$ acting discretely and cocompactly on a
  Euclidean space is the fundamental group of a compact manifold with
  nonnegative sectional curvature.
  \source{petersen:page-0320}
\end{remark}

\begin{corollary}[Corollary 7.3.13]
  \label{cor:pet-ch7-flat-manifold-Kpi1}
  \lean{PetersenLib.flatIfAsphericalNonnegRicci}
  \uses{thm:pet-ch7-structure-theorem}
  Suppose $(M,g)$ is compact with $\Ric\ge0$. If $M$ is a $K(\pi,1)$ (its
  universal cover is contractible), then the universal cover is Euclidean
  space and $(M,g)$ is flat.
  \source{petersen:page-0320,petersen:page-0321}
\end{corollary}
\begin{proof}
  \uses{cor:pet-ch7-flat-manifold-Kpi1}
  By \cref{thm:pet-ch7-structure-theorem}(1), $\tilde M=\mathbb{R}^k\times C$ with
  $C$ compact. This space is contractible only if $C$ is contractible, and
  the only compact contractible manifold is a point; hence $\tilde
  M=\mathbb{R}^k=\mathbb{R}^n$, with the flat metric forced by the product splitting
  degenerating to a single Euclidean factor.
\end{proof}

\begin{corollary}[Corollary 7.3.14]
  \label{cor:pet-ch7-finite-fundamental-group-positive-ricci}
  \lean{PetersenLib.finitePi1IfPositiveRicciSomewhere}
  \uses{thm:pet-ch7-structure-theorem}
  If $(M,g)$ is compact with $\Ric\ge0$ and $\Ric>0$ on some tangent space
  $T_pM$, then $\pi_1(M)$ is finite.
  \source{petersen:page-0321}
\end{corollary}
\begin{proof}
  \uses{cor:pet-ch7-finite-fundamental-group-positive-ricci}
  Positivity of $\Ric$ on an entire tangent space prevents the universal
  cover from splitting off any $\mathbb{R}^k$ factor with $k\ge1$ (the splitting
  theorem would force $\Ric(\partial_t,\partial_t)=0$ along the $\mathbb{R}$-direction
  of any such factor). Hence in \cref{thm:pet-ch7-structure-theorem}(1),
  $k=0$ and $\tilde M=N$ is compact, so $\pi_1(M)$ acts discretely and
  cocompactly on the compact space $\tilde M$, forcing $\pi_1(M)$ finite.
\end{proof}

\begin{corollary}[Corollary 7.3.15]
  \label{cor:pet-ch7-first-betti-number-bound}
  \lean{PetersenLib.firstBettiNumberBoundNonnegRicci}
  \uses{thm:pet-ch7-structure-theorem}
  If $(M,g)$ is compact with $\Ric\ge0$, then $b_1(M)\le\dim M=n$, with
  equality iff $(M,g)$ is a flat torus.
  \source{petersen:page-0321}
\end{corollary}
\begin{proof}
  \uses{cor:pet-ch7-first-betti-number-bound}
  There is a natural surjection $h:\pi_1(M)\to H_1(M,\mathbb{Z})\cong\mathbb{Z}^{b_1}\times
  T$ ($T$ finite abelian) sending loops to cycles. Since the finite normal
  subgroup $G$ of \cref{thm:pet-ch7-structure-theorem}(3) is finite, $h(G)$
  lies in the torsion part $T$, so $h$ descends to a surjection
  $\pi_1(M)/G=\pi_1(M)\cap\Iso(\mathbb{R}^k)\twoheadrightarrow\mathbb{Z}^{b_1}$. The image of
  the finite-index subgroup $\mathbb{Z}^k\subset\pi_1(M)\cap\Iso(\mathbb{R}^k)$ then also
  has finite index in $\mathbb{Z}^{b_1}$ (a finite-index subgroup maps onto a
  finite-index subgroup of the image under a surjection with finite-index
  domain restriction), giving $b_1\le k\le n$.

  If $b_1=n$, then $k=n$ forces $\tilde M=\mathbb{R}^n$ (no compact factor $N$
  remains, i.e.\ $N$ is a point), so $G$ is trivial and $h:\mathbb{Z}^n\to\mathbb{Z}^{b_1}=
  \mathbb{Z}^n$ restricted to the finite-index $\mathbb{Z}^n\subset\pi_1(M)$ must be
  injective (else its image would fail to have finite index in $\mathbb{Z}^n$). So
  $\ker h$ meets the finite-index subgroup $\mathbb{Z}^n\subset\pi_1(M)$ trivially,
  hence $\ker h$ is finite; but any finite-order isometry of $\mathbb{R}^n$ has a
  fixed point, and elements of $\ker h\subset\pi_1(M)$ act freely (as deck
  transformations), so $\ker h$ is trivial. Thus $\pi_1(M)\cong\mathbb{Z}^n\times T$
  with $T$ (now seen to be) trivial, so $M=\mathbb{R}^n/\mathbb{Z}^n$ is a torus (not
  necessarily the standard square torus, since the $\mathbb{Z}^n$-action on
  $\mathbb{R}^n$ need not be the standard one).
\end{proof}

\begin{theorem}[Theorem 7.3.16]
  \label{thm:pet-ch7-homogeneous-splitting}
  \lean{PetersenLib.homogeneousNonnegRicciSplitting}
  \uses{thm:pet-ch7-splitting-theorem,thm:pet-ch7-structure-theorem}
  Let $(M,g)$ be a homogeneous Riemannian manifold with $\Ric\ge0$. Then
  $(M,g)=(N\times\mathbb{R}^k,\,g_N+g_{\mathbb{R}^k})$ with $(N,g_N)$ compact homogeneous.
  \source{petersen:page-0321,petersen:page-0322}
\end{theorem}
\begin{proof}
  \uses{thm:pet-ch7-homogeneous-splitting}
  First split $(M,g)=(N\times\mathbb{R}^k,g_N+g_{\mathbb{R}^k})$
  (\cref{thm:pet-ch7-splitting-theorem}, iterated as in
  \cref{thm:pet-ch7-structure-theorem}(1)) so that $N$ contains no lines. As
  in \cref{thm:pet-ch7-structure-theorem}(2), the isometry group of $M$
  splits as $\Iso(N)\times\Iso(\mathbb{R}^k)$, and since $\Iso(M)$ acts transitively
  and preserves the product structure, $\Iso(N)$ acts transitively on $N$: it
  makes $N$ homogeneous.

  It remains to show any noncompact homogeneous space contains a line (which
  by construction of $N$ then forces $N$ to be compact, as in
  \cref{thm:pet-ch7-structure-theorem}(3)). Let $c:[0,\infty)\to M$ be a
  unit-speed ray (\cref{lem:pet-ch7-ray-line-existence}) and, using
  homogeneity, choose isometries $F_s$ with $F_s(c(s))=c(0)$ for each $s\ge0$.
  Define $c_s:[-s,\infty)\to M$ by $c_s(t)=F_s(c(t+s))$; then $c_s(0)=c(0)$
  and $\dot c_s(0)=DF_s(\dot c(s))$. Since $F_s$ is an isometry taking the
  segment $c|_{[0,s]}$ (from $c(0)$ to $c(s)$, realizing the distance $s$) to
  a segment from $F_s(c(0))$ to $c(0)$ of the same length, and $c_s$ is a
  reparametrized ray by construction, in the limit (extracting a convergent
  subsequence of $\dot c_s(0)$ as $s\to\infty$, using that each $c_s$ is a
  ray on $[-s,\infty)$ so any fixed compact interval $[-a,\infty)$ eventually
  lies in the domain) the extension of $c$ backward along the limiting
  direction remains distance-realizing on every finite interval, producing a
  full line through $c(0)$.
\end{proof}

\section{Further Study}

The interested reader may consult additional treatments of the material in
this chapter in the wider literature: surveys of comparison geometry more
generally; Anderson's article containing further examples of manifolds with
nonnegative Ricci curvature; the papers constructing examples with almost
maximal diameter (\cref{ex:pet-ch7-anderson-cp2-example,ex:pet-ch7-otsu-example});
the original paper of Cheeger and Gromoll on the splitting theorem
(\cref{thm:pet-ch7-splitting-theorem}) together with a later elementary proof
of it; and the articles of Colding, Perel'man, and Zhu tracing the subsequent
development of the subject.
\source{petersen:page-0322}

\section{Exercises}

\begin{remark}[Exercise 7.5.1]
  \label{rem:pet-ch7-ex-1}
  \lean{PetersenLib.exercise7_5_1}
  \uses{prop:pet-ch7-volume-form-laplacian-trace,def:pet-ch7-polar-density}
  With notation as in §7.1.1 and using $\vol=\lambda\,dr\wedge\vol_{n-1}$, show
  that $\mu=\lambda^{1/(n-1)}$ satisfies
  \[
    \partial_r^2\mu\le-\frac{\mu}{n-1}\Ric(\partial_r,\partial_r),\qquad
    \mu(0,\theta)=0,\qquad \lim_{r\to0}\partial_r\mu(r,\theta)=1,
  \]
  and that this yields the desired volume-form estimates.
  \source{petersen:page-0322}
\end{remark}

\begin{remark}[Exercise 7.5.2 — Calabi and Yau]
  \label{rem:pet-ch7-ex-2}
  \lean{PetersenLib.exercise7_5_2}
  \uses{thm:pet-ch7-structure-theorem}
  Let $(M,g)$ be complete noncompact with $\Ric\ge0$, $p\in M$.
  \begin{enumerate}
    \item[(1)] Show that for each $R>1$ there is $x\in M$ with
    \[
      \vol B(p,1)\le\vol B(x,R+1)-\vol B(x,R-1)
      \le\frac{(R+1)^n-(R-1)^n}{(R+1)^n}\vol B(p,2R).
    \]
    \item[(2)] Show there is $C>0$ with $\vol B(p,R)\ge CR$.
  \end{enumerate}
  \source{petersen:page-0322}
\end{remark}

\begin{remark}[Exercise 7.5.3]
  \label{rem:pet-ch7-ex-3}
  \lean{PetersenLib.exercise7_5_3}
  \uses{def:pet-ch7-support-function}
  Let $f:I\to\mathbb{R}$ be continuous on an interval $I\subset\mathbb{R}$. Show that (1)
  $f$ is convex, (2) $f$ has a linear support function from below of the form
  $a(x-x_0)+f(x_0)$ at every $x_0\in I$, and (3) $f''\ge0$ in the support
  sense at every $x_0\in I$, are equivalent.
  \source{petersen:page-0322,petersen:page-0323}
\end{remark}

\begin{remark}[Exercise 7.5.4]
  \label{rem:pet-ch7-ex-4}
  \lean{PetersenLib.exercise7_5_4}
  \uses{thm:pet-ch7-poincare-sobolev}
  Show that on a compact Riemannian manifold there is no finite $S(s)$ with
  $\|f-f_M\|_\infty\le S\|df\|_1$ for $1<s<\dim M$.
  \source{petersen:page-0323}
\end{remark}

\begin{remark}[Exercise 7.5.5 — Basic Covering Lemma]
  \label{rem:pet-ch7-ex-5}
  \lean{PetersenLib.exercise7_5_5}
  \uses{thm:pet-ch7-maximal-function-theorem}
  Given a separable metric space $(X,d)$ and a bounded positive function
  $R:X\to(0,D]$, show there is a countable $A\subset X$ such that the balls
  $B(p,R(p))$, $p\in A$, are pairwise disjoint and $X=\bigcup_{p\in
  A}B(p,5R(p))$. (Hint: choose the points successively so that
  $R(p_{k+1})\ge\tfrac12\sup_{p\in X\setminus\bigcup_iB(p_i,2R(p_i))}R(p)$.)
  \source{petersen:page-0323}
\end{remark}

\begin{remark}[Exercise 7.5.6]
  \label{rem:pet-ch7-ex-6}
  \lean{PetersenLib.exercise7_5_6}
  \uses{thm:pet-ch7-cheng-rigidity}
  Assume $r(x)=|xp|$ is smooth on $B(p,R)$. Show that if
  $\Hess r=(\snk'(r)/\snk(r))g_r$ in polar coordinates, then all sectional
  curvatures on $B(p,R)$ equal $k$.
  \source{petersen:page-0323}
\end{remark}

\begin{remark}[Exercise 7.5.7]
  \label{rem:pet-ch7-ex-7}
  \lean{PetersenLib.exercise7_5_7}
  \uses{thm:pet-ch7-poincare-sobolev}
  Construct convex surfaces in $\mathbb{R}^3$ by capping off cylinders
  $[-R,R]\times S^1$ to show the Sobolev--Poincaré constants increase as $R$
  increases. (Hint: consider test functions constant except on $[-1,1]\times
  S^1$.)
  \source{petersen:page-0323}
\end{remark}

\begin{remark}[Exercise 7.5.8]
  \label{rem:pet-ch7-ex-8}
  \lean{PetersenLib.exercise7_5_8}
  \uses{lem:pet-ch7-absolute-volume-comparison}
  Show that if $(M,g)$ has $\Ric\ge(n-1)k$ and $\vol B(p,R)=v(n,k,R)$ for some
  $p\in M$, then the metric has constant curvature $k$ on $B(p,R)$.
  \source{petersen:page-0323}
\end{remark}

\begin{remark}[Exercise 7.5.9]
  \label{rem:pet-ch7-ex-9}
  \lean{PetersenLib.exercise7_5_9}
  \uses{lem:pet-ch7-laplacian-comparison-calabi}
  Let $X$ be a vector field and $F_t(p)=\exp_p(tX|_p)$.
  \begin{enumerate}
    \item[(1)] For $v\in T_pM$, show $J(t)=DF_t(v)$ is a Jacobi field along
      $c(t)=\exp(tX)$ with $J(0)=v$, $\dot J(0)=\nabla_vX$.
    \item[(2)] For an orthonormal basis $e_i$ of $T_pM$ and $J_i(t)=DF_t(e_i)$,
      show $(\det[DF_t])^2=\det[g(J_i(t),J_j(t))]$.
    \item[(3)] Show that as long as $\det(DF_t)\ne0$,
      \[
        \frac{d^2(\det(DF_t))^{1/n}}{dt^2}
        \le-\frac{(\det(DF_t))^{1/n}}{n}\Ric(\dot c,\dot c).
      \]
      (Hint: use $(\operatorname{tr}A)^2\le n\operatorname{tr}(A^*A)$ for
      $n\times n$ matrices $A$.)
  \end{enumerate}
  \source{petersen:page-0323}
\end{remark}

\begin{remark}[Exercise 7.5.10]
  \label{rem:pet-ch7-ex-10}
  \lean{PetersenLib.exercise7_5_10}
  \uses{lem:pet-ch7-relative-volume-comparison}
  Show that a complete manifold $(M,g)$ with $\Ric\ge0$ and
  $\lim_{r\to\infty}\vol B(p,r)/(\omega_nr^n)=1$ for some $p\in M$ must be
  isometric to Euclidean space.
  \source{petersen:page-0324}
\end{remark}

\begin{remark}[Exercise 7.5.11]
  \label{rem:pet-ch7-ex-11}
  \lean{PetersenLib.exercise7_5_11}
  \uses{lem:pet-ch7-hessian-touching-comparison}
  Show that any function on an $n$-dimensional Riemannian manifold satisfies
  $|\Hess u|^2\ge\tfrac1n|\Delta u|^2$, with equality iff
  $\Hess u=\tfrac{\Delta u}{n}g$. What can be said about $M$ when
  $\Hess u=\tfrac{\Delta u}{n}g$?
  \source{petersen:page-0324}
\end{remark}

\begin{remark}[Exercise 7.5.12]
  \label{rem:pet-ch7-ex-12}
  \lean{PetersenLib.exercise7_5_12}
  \uses{def:pet-ch7-linear-harmonic-function}
  Show that if $u,v:M\to\mathbb{R}$ are compactly supported and smooth on open dense
  sets, then
  \[
    \int u\,\Delta v\,\vol=\int v\,\Delta u\,\vol
    =-\int g(du,dv)\,\vol=-\int g(\nabla u,\nabla v)\,\vol.
  \]
  \source{petersen:page-0324}
\end{remark}

\begin{remark}[Exercise 7.5.13]
  \label{rem:pet-ch7-ex-13}
  \lean{PetersenLib.exercise7_5_13}
  \uses{thm:pet-ch7-strong-maximum-principle}
  Show that if $\Delta u=\lambda u$ on a closed Riemannian manifold, then
  $\lambda\le0$, and $\lambda=0$ forces $u$ constant.
  \source{petersen:page-0324}
\end{remark}

\begin{remark}[Exercise 7.5.14]
  \label{rem:pet-ch7-ex-14}
  \lean{PetersenLib.exercise7_5_14}
  \uses{thm:pet-ch7-cheng-rigidity}
  Show that $u_k=\cos(\sqrt kr)$ on $S^n_k=S^n(1/\sqrt k)$ satisfies $\Delta
  u_k=-nk\,u_k$ and $\int u_k\,\vol=0$.
  \source{petersen:page-0324}
\end{remark}

\begin{remark}[Exercise 7.5.15 — Lichnerowicz]
  \label{rem:pet-ch7-ex-15}
  \lean{PetersenLib.exercise7_5_15}
  \uses{thm:pet-ch7-cheng-rigidity}
  Let $(M^n,g)$ be closed with $\Ric\ge(n-1)k>0$. Use the Bochner formula to
  show that every $u$ with $\Delta u=-\lambda u$, $\lambda>0$, satisfies
  $\lambda\ge nk$. Deduce, via the spectral theorem for $\Delta$, that every
  $u$ with $\int u\,\vol=0$ satisfies the Poincaré inequality
  $\int u^2\,\vol\le\tfrac1{nk}\int|du|^2\,\vol$.
  \source{petersen:page-0324}
\end{remark}

\begin{remark}[Exercise 7.5.16 — Obata]
  \label{rem:pet-ch7-ex-16}
  \lean{PetersenLib.exercise7_5_16}
  \uses{rem:pet-ch7-ex-15,thm:pet-ch7-cheng-rigidity}
  Let $(M^n,g)$ be closed with $\Ric\ge(n-1)k>0$. Using the Bochner formula as
  in \cref{rem:pet-ch7-ex-15}, show that if $\Delta u=-nk\,u$ for some
  function $u$, then $\Hess u=-ku\,g$. Conclude $(M^n,g)=S^n_k$.
  \source{petersen:page-0324}
\end{remark}

\begin{remark}[Exercise 7.5.17 — P. Li and Schoen]
  \label{rem:pet-ch7-ex-17}
  \lean{PetersenLib.exercise7_5_17}
  \uses{thm:pet-ch7-poincare-sobolev}
  Let $(M,g)$ be complete with $\Ric\ge-(n-1)k^2$, $k\ge0$; $p\in M$; $R>0$
  with $\partial B(p,2R)\ne\emptyset$; $q\in\partial B(p,2R)$;
  $r(x)=|xq|$.
  \begin{enumerate}
    \item[(1)] Show $\Delta r\le(n-1)(R^{-1}+k)$ on $B(p,R)$.
    \item[(2)] Let $f(x)=a\exp(-ar(x))$, $a>0$. Show
      $\Delta f\ge a\exp(-3aR)\big(a-(n-1)(R^{-1}+k)\big)$.
    \item[(3)] Let $u\ge0$ be smooth with compact support in $B(p,R)$, and
      choose $a=n(R^{-1}+k)$. Using
      $\int_{B(p,R)}u\,\Delta f\,\vol=-\int_{B(p,R)}g(du,df)\,\vol$, show
      $\int_{B(p,R)}u\,\vol\le C\int_{B(p,R)}|du|\,\vol$, where
      $C=\dfrac{R}{1+kR}\exp\big(2n(1+kR)\big)$.
    \item[(4)] Extend the inequality of (3) to all smooth $u$ with compact
      support in $B(p,R)$.
    \item[(5)] For $s\ge1$ and $u$ with compact support in $B(p,R)$, show
      $\int_{B(p,R)}|u|^s\,\vol\le(sC)^s\int_{B(p,R)}|du|\,\vol$.
  \end{enumerate}
  \source{petersen:page-0324,petersen:page-0325}
\end{remark}

\begin{remark}[Exercise 7.5.18 — Cheeger]
  \label{rem:pet-ch7-ex-18}
  \lean{PetersenLib.exercise7_5_18}
  \uses{lem:pet-ch7-relative-volume-comparison}
  Suppose $(M^n,g)$ has $\Ric\ge(n-1)k$.
  \begin{enumerate}
    \item[(1)] For $p_1,\dots,p_k\in M$, show
      $r\mapsto\vol\big(\bigcup_{i=1}^kB(p_i,r)\big)/v(n,k,r)$ is
      nonincreasing and converges to $k$ as $r\to0$.
    \item[(2)] For $A\subset M$, show
      $r\mapsto\vol\big(\bigcup_{p\in A}B(p,r)\big)/v(n,k,r)$ is
      nonincreasing, by applying (1) to a dense finite subset of $A$.
  \end{enumerate}
  \source{petersen:page-0325}
\end{remark}

\begin{remark}[Exercise 7.5.19]
  \label{rem:pet-ch7-ex-19}
  \lean{PetersenLib.exercise7_5_19}
  \uses{lem:pet-ch7-absolute-volume-comparison}
  For $p\in M$ and $\Gamma\subset T_pM$ a set of unit vectors, define the
  polar cone $B^\Gamma(p,R)=\{(t,\theta)\in M\mid t\le R,\ \theta\in\Gamma\}$.
  If $\Ric M\ge(n-1)k$, show
  \[
    \vol B^\Gamma(p,R)\le\vol\Gamma\cdot\int_0^R(\snk_k(t))^{n-1}\,dt.
  \]
  \source{petersen:page-0326}
\end{remark}

\begin{remark}[Exercise 7.5.20]
  \label{rem:pet-ch7-ex-20}
  \lean{PetersenLib.exercise7_5_20}
  \uses{thm:pet-ch7-structure-theorem}
  Let $G$ be a compact connected Lie group with a biinvariant metric and
  $\Ric\ge0$.
  \begin{enumerate}
    \item[(1)] If $G$ has finite center, show $\pi_1(G)$ is finite.
    \item[(2)] Show a finite cover of $G$ has the form $G'\times T^k$ with
      $G'$ compact simply connected and $T^k$ a torus.
    \item[(3)] If $\pi_1(G)$ is finite, show the center of $G$ is finite.
  \end{enumerate}
  \source{petersen:page-0326}
\end{remark}

\begin{remark}[Exercise 7.5.21]
  \label{rem:pet-ch7-ex-21}
  \lean{PetersenLib.exercise7_5_21}
  \uses{thm:pet-ch7-structure-theorem}
  Let $(M,g)$ be $n$-dimensional and isometric to Euclidean space outside a
  compact $K\subset M$, i.e.\ $M-K$ is isometric to $\mathbb{R}^n-C$ for some
  compact $C\subset\mathbb{R}^n$. If $\Ric\ge0$, show $M=\mathbb{R}^n$.
  \source{petersen:page-0326}
\end{remark}

\begin{remark}[Exercise 7.5.22]
  \label{rem:pet-ch7-ex-22}
  \lean{PetersenLib.exercise7_5_22}
  \uses{thm:pet-ch7-cheng-rigidity}
  Show that if $\Ric\ge n-1$ then $\diam\le\pi$, by showing that if
  $|pq|>\pi$, then $e_{p,q}(x)=|px|+|xq|-|pq|$ has negative Laplacian at a
  local minimum.
  \source{petersen:page-0326}
\end{remark}
